\section{Case Study II: Exact Analysis of an Educational Spy Detector}
\label{sec:case-study-ii}

This case study began with something visually simple: the baseline circuit had
one amplitude pattern, the ordinary spy changed it, and the advanced circuit
seemed to put the pattern back. That observation was interesting, but it was not
yet a result. A picture can hide phase, wire order, discarded registers, and
correlation. The useful question is therefore not merely whether two amplitude
displays look alike. It is:

\begin{quote}
What, exactly, is the same at the end of the circuit, and what has quietly
changed somewhere that the displayed histogram does not show?
\end{quote}

The source example is the ``Quantum Spy Hunter'' circuit from
\emph{Programming Quantum Computers} \cite{johnston2019programming}. The
maintained errata corrects the unconditional detection probability for the
simple intercept-and-resend construction from $25\%$ to $12.5\%$
\cite{oreilly2026errata}. The author later supplied the machine-readable Quirk
column arrays and complete five-wire output amplitudes for three circuits:

\begin{enumerate}
  \item a baseline with no spy;
  \item an ordinary intercept-and-resend construction; and
  \item an advanced source-access construction.
\end{enumerate}

The identity blocks labeled \texttt{inputA3}, \texttt{inputR2}, and
\texttt{inputB3} are notes rather than operations. Likewise, the probability,
Bloch, density, and amplitude displays observe the simulation but do not change
the state. The circuit figures below omit those non-operational widgets and
replace them with separate amplitude-density figures. Every actual gate in the
supplied column arrays is retained.

\subsection{Reading the circuit literally}

There are three details worth stating before doing any algebra.

First, the supplied reconstruction contained a repeated measurement checkpoint
on Alice's selector wires. No operation changes those selector wires between
the two measurements, so the second checkpoint is redundant in the ideal
projective model. Revision 0.6 and the corrected repository artifact retain one
Alice checkpoint immediately after the two initial Hadamards. This correction
clarifies the intended semantic flow without changing the branch values used by
later classically controlled operations.

Second, the controls on $q_1$ and $q_3$ act on Hadamard gates. They are
controlled-$H$ operations, not an unconnected control marker followed by an
independent Hadamard. The vertical connections are explicit in all three
figures.

Third, the two routing operations are ordered rather than simultaneous. The
ordinary and advanced circuits apply
\begin{equation}
  \operatorname{SWAP}(q_3,q_4)
  \quad\text{followed by}\quad
  \operatorname{SWAP}(q_2,q_4).
  \label{eq:caseii-swap-order}
\end{equation}
They are drawn in separate columns with a small visual spacer. The spacer is an
identity in time, not a new operation; its only purpose is to keep the two SWAP
lines from becoming one unreadable knot.

The meter symbols should also be read with Quirk's simulation semantics in
mind. Later controls conditioned on a measured wire are classical in the
operational interpretation, while Quirk can retain a deferred-measurement
purification internally in order to produce a complete amplitude export. The
paper uses the gate order exactly as supplied and uses the exported amplitudes,
rather than an independently guessed collapse history, for the state analysis.


\begin{figure}[tbp]
\centering
\begin{adjustbox}{max width=\textwidth}
\begin{quantikz}[row sep=0.34cm,column sep=0.30cm]
\lstick{$q_0$} & \gate{H} & \meter{} & \ctrl{2} & \qw      & \slice[style={black,dashed},label style={black}]{Alice checkpoint} \qw & \qw      & \qw      & \qw       & \qw \\
\lstick{$q_1$} & \gate{H} & \meter{} & \qw      & \ctrl{1} & \qw & \qw      & \qw      & \qw       & \qw \\
\lstick{$q_2$} & \qw      & \qw      & \targ{}  & \gate{H} & \qw & \qw      & \qw      & \gate{H}  & \meter{} \\
\lstick{$q_3$} & \qw      & \qw      & \qw      & \qw      & \qw & \gate{H} & \meter{} & \ctrl{-1} & \qw \\
\lstick{$q_4$} & \qw      & \qw      & \qw      & \qw      & \qw & \qw      & \qw      & \qw       & \qw
\end{quantikz}
\end{adjustbox}
\caption{Corrected baseline circuit. One Alice measurement checkpoint remains
after the initial Hadamards; the redundant consecutive checkpoint has been
removed. The control on $q_1$ acts on a Hadamard target on $q_2$; later, the
measured value on $q_3$ controls a second Hadamard on $q_2$. Display widgets and
identity annotation blocks are omitted.}
\label{fig:caseii-baseline-circuit}
\end{figure}

\begin{figure}[tbp]
\centering
\begin{adjustbox}{max width=\textwidth}
\begin{quantikz}[row sep=0.34cm,column sep=0.31cm]
\lstick{$q_0$} & \gate{H} & \meter{} & \ctrl{2} & \qw      & \slice[style={black,dashed},label style={black}]{Alice} \qw & \qw      & \qw      & \qw      & \qw      & \qw      & \slice[style={black,dashed},label style={black}]{routing} \qw & \qw      & \qw      & \qw       & \qw \\
\lstick{$q_1$} & \gate{H} & \meter{} & \qw      & \ctrl{1} & \qw & \qw      & \qw      & \qw      & \qw      & \qw      & \qw & \qw      & \qw      & \qw       & \qw \\
\lstick{$q_2$} & \qw      & \qw      & \targ{}  & \gate{H} & \qw & \gate{H} & \meter{} & \qw      & \qw      & \swap{2} & \qw & \qw      & \qw      & \gate{H}  & \meter{} \\
\lstick{$q_3$} & \qw      & \qw      & \qw      & \qw      & \qw & \gate{H} & \gate{X} & \swap{1} & \qw      & \qw      & \qw & \gate{H} & \meter{} & \ctrl{-1} & \qw \\
\lstick{$q_4$} & \qw      & \qw      & \qw      & \qw      & \qw & \qw      & \qw      & \targX{} & \qw      & \targX{} & \qw & \qw      & \qw      & \qw       & \qw
\end{quantikz}
\end{adjustbox}
\caption{Ordinary intercept--resend circuit. The routing is the ordered pair
$\operatorname{SWAP}(q_3,q_4)$ followed by
$\operatorname{SWAP}(q_2,q_4)$. The operations occupy separate columns, with a
visual spacer between them, so the SWAP lines remain distinct. The $H$ on
$q_2$ controlled by $q_3$ is drawn explicitly near the end of the circuit.}
\label{fig:caseii-intercept-circuit}
\end{figure}

\begin{figure}[tbp]
\centering
\begin{adjustbox}{max width=\textwidth}
\begin{quantikz}[row sep=0.34cm,column sep=0.27cm]
\lstick{$q_0$} & \gate{H} & \meter{} & \ctrl{3} & \qw      & \ctrl{2} & \qw      & \slice[style={black,dashed},label style={black}]{Alice} \qw & \qw      & \qw      & \qw      & \qw      & \qw      & \slice[style={black,dashed},label style={black}]{routing} \qw & \qw      & \qw      & \qw       & \qw \\
\lstick{$q_1$} & \gate{H} & \meter{} & \qw      & \ctrl{2} & \qw      & \ctrl{1} & \qw & \qw      & \qw      & \qw      & \qw      & \qw      & \qw & \qw      & \qw      & \qw       & \qw \\
\lstick{$q_2$} & \qw      & \qw      & \qw      & \qw      & \targ{}  & \gate{H} & \qw & \gate{H} & \meter{} & \qw      & \qw      & \swap{2} & \qw & \qw      & \qw      & \gate{H}  & \meter{} \\
\lstick{$q_3$} & \qw      & \qw      & \targ{}  & \gate{H} & \qw      & \qw      & \qw & \qw      & \qw      & \swap{1} & \qw      & \qw      & \qw & \gate{H} & \meter{} & \ctrl{-1} & \qw \\
\lstick{$q_4$} & \qw      & \qw      & \qw      & \qw      & \qw      & \qw      & \qw & \qw      & \qw      & \targX{} & \qw      & \targX{} & \qw & \qw      & \qw      & \qw       & \qw
\end{quantikz}
\end{adjustbox}
\caption{Corrected advanced source-access circuit. Alice's measured selectors
control preparation operations on $q_3$ and then a second preparation on $q_2$;
the redundant repeated selector measurement has been removed. The ordered SWAP
pair remains in separate columns.}
\label{fig:caseii-advanced-circuit}
\end{figure}


For reference, the operational columns represented by the three diagrams are
summarized in \cref{tab:caseii-gate-order}. Display-only columns and identity
labels are not listed.

\begin{table}[ht]
\centering
\caption{Gate order reconstructed from the supplied Quirk column arrays.}
\label{tab:caseii-gate-order}
\begin{tabular}{@{}p{0.16\textwidth}p{0.75\textwidth}@{}}
\toprule
Variant & Ordered operations \\
\midrule
Baseline &
$H(q_0),H(q_1)$; $M(q_0),M(q_1)$;
$\cnot(q_0,q_2)$; $C_H(q_1,q_2)$;
$H(q_3)$; $M(q_3)$; $C_H(q_3,q_2)$; $M(q_2)$.\\[0.4em]
Intercept--resend &
$H(q_0),H(q_1)$; $M(q_0),M(q_1)$;
$\cnot(q_0,q_2)$; $C_H(q_1,q_2)$;
$H(q_2),H(q_3)$; $M(q_2),X(q_3)$;
$\operatorname{SWAP}(q_3,q_4)$;
$\operatorname{SWAP}(q_2,q_4)$;
$H(q_3)$; $M(q_3)$; $C_H(q_3,q_2)$; $M(q_2)$.\\[0.4em]
Advanced &
$H(q_0),H(q_1)$; $M(q_0),M(q_1)$;
$\cnot(q_0,q_3)$; $C_H(q_1,q_3)$;
$\cnot(q_0,q_2)$; $C_H(q_1,q_2)$;
$H(q_2)$; $M(q_2)$;
$\operatorname{SWAP}(q_3,q_4)$;
$\operatorname{SWAP}(q_2,q_4)$;
$H(q_3)$; $M(q_3)$; $C_H(q_3,q_2)$; $M(q_2)$.\\
\bottomrule
\end{tabular}
\end{table}

\subsection{The baseline detection rule}

Let Alice choose a value bit $v$ and a basis bit $b$. The corresponding BB84
state is
\begin{equation}
  \ket{\psi_{v,b}}=H^bX^v\ket{0}.
  \label{eq:caseii-bb84-state}
\end{equation}
The four cases are $\ket{0}$, $\ket{1}$, $\ket{+}$, and $\ket{-}$. Bob chooses
his own basis bit $c$, applies $H^c$, and measures a value $y$. If Bob chose the
same basis as Alice, then
\begin{equation}
  H^c\ket{\psi_{v,b}}
  =H^bH^bX^v\ket{0}
  =\ket{v},
  \qquad c=b.
\end{equation}
The ideal detection event is therefore
\begin{equation}
  D=\{c=b\text{ and }y\neq v\}.
  \label{eq:caseii-detection-event}
\end{equation}
Without a spy, $\Prb(D)=0$.

\subsection{Why ordinary intercept-and-resend is visible}

An ordinary attacker selects a random basis $e$, measures the traveling qubit,
and sends a replacement prepared from the observed result. Three independent
things must happen for the educational circuit to catch the attacker on a
given unsifted transmission:

\begin{enumerate}
  \item Bob must choose Alice's basis, which occurs with probability $1/2$;
  \item the attacker must choose the wrong basis, also with probability $1/2$;
  \item Bob must receive the wrong value after that wrong-basis intervention,
  which occurs with probability $1/2$.
\end{enumerate}

\begin{proposition}[Detection probability]
For uniformly random independent $b$, $c$, and $e$,
\begin{equation}
  \Prb(D)=\frac12\cdot\frac12\cdot\frac12
  =\frac18=12.5\%.
  \label{eq:caseii-detection-probability}
\end{equation}
Conditioned on Alice and Bob selecting the same basis, the induced error rate is
$25\%$.
\end{proposition}

This is often summarized by saying that measurement disturbs the state. That
summary is true, but it is too vague to carry a security proof. The actual
reason is that the BB84 preparation family contains nonorthogonal states. A
measurement that distinguishes one basis cannot preserve the conjugate basis in
every case. The circuit makes that geometry visible as an altered output
distribution \cite{fuchs1996information}.

\subsection{The advanced construction and the source-access interpretation}

The advanced circuit uses Alice's selector wires twice. They first control an
$X$ and a controlled-$H$ preparation on $q_3$, and later control the same type
of preparation on $q_2$. At an abstract level, this is consistent with
executing a known preparation family on two initialized targets:
\begin{equation}
  \ket{v,b}\ket{0}_S\ket{0}_E
  \longmapsto
  \ket{v,b}\ket{\psi_{v,b}}_S\ket{\psi_{v,b}}_E.
  \label{eq:caseii-prepare-twice}
\end{equation}

This is not universal cloning. The circuit is not handed an arbitrary unknown
$\ket{\psi}$ and asked to discover how to copy it. It is given access to the
controls that define a finite preparation procedure and runs that procedure
more than once. There is no contradiction with the no-cloning theorem
\cite{wootters1982cloning}.

That distinction also changes the security model. A channel-only spy sees a
traveling quantum state. A source-access spy can reach the value and basis
controls, insert another target into the preparation path, or obtain a
correlated emission from the source. The advanced circuit should therefore be
read as a source-compromise or privileged-compilation example, not as a proof
that standard BB84 is broken under its ordinary assumptions
\cite{tamaki2014loss,jiang2024sidechannel}.

The abstract interpretation is useful, but the complete export lets us say
something stronger and more exact without guessing the semantic name of every
wire.

\subsection{Semantic wire roles through the routing network}

The two ordered SWAPs implement the three-wire permutation
\begin{equation}
  (q_2,q_3,q_4)\longmapsto(q_3,q_4,q_2).
  \label{eq:caseii-semantic-permutation}
\end{equation}
The physical wires do not move, but their logical responsibilities rotate.
\Cref{tab:caseii-semantic-routing} records the interpretation used by the
fixed-input implementation.

\begin{table}[ht]
\centering
\caption{Semantic roles immediately before and after the two SWAPs.}
\label{tab:caseii-semantic-routing}
\begin{tabularx}{\textwidth}{@{}cXX@{}}
\toprule
Physical wire & Before routing & After routing \\
\midrule
$q_2$ & Original signal after spy-side basis rotation and measurement &
Receiver-visible signal, inherited from the separately prepared $q_3$ branch \\
$q_3$ & Separately prepared replacement or duplicate signal & Clean workspace,
inherited from $q_4$, used for the receiver's basis selection \\
$q_4$ & Clean routing workspace initialized to $\ket{0}$ & Retained spy-side
branch, inherited from the former $q_2$ state \\
\bottomrule
\end{tabularx}
\end{table}

The fifth physical wire is therefore the natural candidate latch. The security
property follows the logical state through the permutation rather than attaching
permanent meaning to a physical wire number.

\subsection{Exact reduced-state reconstruction}

Let $c^{(x)}_{e,j}$ be the complete output amplitude for circuit variant $x$.
The index $e\in\{0,1\}$ labels the two 16-amplitude blocks associated with the
fifth exported wire, and $j\in\{0,\ldots,15\}$ labels the four-wire register
covered by the final amplitude display. After normalizing the supplied arrays,
the reduced state visible to that four-wire display is
\begin{equation}
  \rho_x[j,k]
  =\sum_{e=0}^{1}c^{(x)}_{e,j}c^{(x)*}_{e,k}.
  \label{eq:caseii-partial-trace}
\end{equation}
This is simply a partial trace written directly in the block ordering used by
the export.

The resulting computational-basis probability densities are shown in
\cref{fig:caseii-probability-densities}. They are not screenshots; the bars are
generated in LaTeX from the complete amplitude arrays.

\begin{figure}[tbp]
\centering
\begin{tikzpicture}
\begin{groupplot}[
group style={group size=1 by 3,vertical sep=0.82cm},
width=0.92\textwidth,height=0.185\textwidth,
xmin=-0.7,xmax=15.7,ymin=0,ymax=0.14,
ytick={0,0.0625,0.125},yticklabels={$0$,$1/16$,$1/8$},
ylabel={probability},
axis line style={black!65},tick style={black!65},
title style={font=\small\bfseries},
grid=major,grid style={black!10},
]
\nextgroupplot[title={Baseline},xtick=\empty]
\addplot[ybar,bar width=4.2pt,fill=BaselineBlue!78,draw=BaselineBlue!95!black] coordinates {(0,0.125000002225) (1,0) (2,0.0624999988874) (3,0.0624999988874) (4,0) (5,0.125000002225) (6,0.0624999988874) (7,0.0624999988874) (8,0.0624999988874) (9,0.0624999988874) (10,0.125000002225) (11,0) (12,0.0624999988874) (13,0.0624999988874) (14,0) (15,0.125000002225)};
\nextgroupplot[title={Intercept--resend},xtick=\empty]
\addplot[ybar,bar width=4.2pt,fill=InterceptOrange!78,draw=InterceptOrange!95!black] coordinates {(0,0.0625000011126) (1,0.0625000011126) (2,0.0624999988874) (3,0.0624999988874) (4,0.0625000011126) (5,0.0625000011126) (6,0.0624999988874) (7,0.0624999988874) (8,0.124999997775) (9,0.124999997775) (10,0.125000002225) (11,0.125000002225) (12,0) (13,0) (14,0) (15,0)};
\nextgroupplot[title={Advanced source access},xtick={0,1,...,15},xticklabels={$0000$,$0001$,$0010$,$0011$,$0100$,$0101$,$0110$,$0111$,$1000$,$1001$,$1010$,$1011$,$1100$,$1101$,$1110$,$1111$},x tick label style={rotate=55,anchor=east,font=\scriptsize},xlabel={display-register basis state}]
\addplot[ybar,bar width=4.2pt,fill=AdvancedTeal!78,draw=AdvancedTeal!95!black] coordinates {(0,0.124999997775) (1,0) (2,0.0624999988874) (3,0.0624999988874) (4,0) (5,0.124999997775) (6,0.0624999988874) (7,0.0624999988874) (8,0.0625000011126) (9,0.0625000011126) (10,0.125000002225) (11,0) (12,0.0625000011126) (13,0.0625000011126) (14,0) (15,0.125000002225)};
\end{groupplot}
\end{tikzpicture}
\caption{Computational-basis probability densities on the four-wire display
register after tracing out the fifth exported wire. The baseline and advanced
panels coincide component by component. The ordinary intercept--resend panel
does not.}
\label{fig:caseii-probability-densities}
\end{figure}

The diagonal vectors are
\begin{align}
\operatorname{diag}(\rho_0)
&=\operatorname{diag}(\rho_A)\nonumber\\
&=\frac1{16}(2,0,1,1,0,2,1,1,1,1,2,0,1,1,0,2),
\label{eq:caseii-matching-diagonals}\\
\operatorname{diag}(\rho_I)
&=\frac1{16}(1,1,1,1,1,1,1,1,2,2,2,2,0,0,0,0),
\label{eq:caseii-intercept-diagonal}
\end{align}
where $0$, $I$, and $A$ denote baseline, intercept--resend, and advanced.

\begin{proposition}[Exact histogram equivalence]
For the supplied exports,
\begin{equation}
  \operatorname{diag}(\rho_A)=\operatorname{diag}(\rho_0)
\end{equation}
up to floating-point export error. The total-variation distance between the two
probability vectors is $8.90\times10^{-9}$.
\end{proposition}

\begin{proof}
For each basis index $j$, Eq.~\eqref{eq:caseii-partial-trace} gives the observed
probability as
\begin{equation}
  p_x(j)=|c^{(x)}_{0,j}|^2+|c^{(x)}_{1,j}|^2.
\end{equation}
Evaluating these sums componentwise for the baseline and advanced exports gives
Eq.~\eqref{eq:caseii-matching-diagonals}. The remaining discrepancy is at the
scale of the floating-point amplitudes supplied by the simulator.
\end{proof}

This is the first important result. A verifier that only repeats the final
computational-basis measurement on those four wires cannot distinguish the
advanced circuit from the baseline. The circuit has restored every answer that
this particular test asks for.

\subsection{The state is not restored}

The interesting part is that matching every displayed probability is still not
the same as matching the quantum state. Probabilities occupy the diagonal of a
density matrix. Relative phase and coherence live in the off-diagonal entries.
The advanced circuit restores the diagonal and changes the rest.

\begin{table}[ht]
\centering
\caption{Exact reduced-state diagnostics calculated from the supplied
five-wire amplitude arrays. Fidelity and trace distance use the pure baseline
state as the reference.}
\label{tab:caseii-density-metrics}
\begin{tabular}{@{}lcccccc@{}}
\toprule
Variant & Quirk coherence & Purity & Entropy & Fidelity & Trace distance & $Z$-TVD\\
\midrule
Baseline & $1.0000$ & $1$ & $0$ & $1$ & $0$ & $0$\\
Intercept--resend & $0.5000$ & $1/2$ & $1$ & $1/4$ & $0.8090$ & $0.3750$\\
Advanced & $0.5000$ & $1/2$ & $1$ & $0.4268$ & $0.6184$ & $8.90\times10^{-9}$\\
\bottomrule
\end{tabular}
\end{table}

For the advanced export,
\begin{equation}
  F(\rho_0,\rho_A)
  =\frac{2+\sqrt2}{8}
  \approx0.4268,
  \label{eq:caseii-fidelity}
\end{equation}
and
\begin{equation}
  D(\rho_0,\rho_A)
  =\frac12\norm{\rho_0-\rho_A}_1
  \approx0.6184.
  \label{eq:caseii-trace-distance}
\end{equation}
With equal prior probabilities, the Helstrom measurement could distinguish the
baseline and advanced reduced states with success probability
\begin{equation}
  P_{\mathrm{distinguish}}
  =\frac{1+D(\rho_0,\rho_A)}{2}
  \approx0.8092.
  \label{eq:caseii-helstrom}
\end{equation}

The reduced density matrices are shown in
\cref{fig:caseii-density-matrices}. Blue cells are positive real entries,
orange cells are negative, and stronger color indicates larger magnitude. The
advanced matrix has exactly the baseline diagonal while its off-diagonal
pattern is visibly different.

\begin{figure}[tbp]
\centering
\begin{tikzpicture}[x=0.255cm,y=0.255cm,font=\scriptsize]
\begin{scope}[xshift=0.0cm]
\node[font=\small\bfseries] at (7.5,1.3) {Baseline $\rho_0$};
\filldraw[fill=BaselineBlue!88,draw=BaselineBlue!45!black,line width=0.08pt] (0,0) rectangle ++(1,-1);
\filldraw[fill=white,draw=black!15,line width=0.08pt] (1,0) rectangle ++(1,-1);
\filldraw[fill=BaselineBlue!62,draw=BaselineBlue!45!black,line width=0.08pt] (2,0) rectangle ++(1,-1);
\filldraw[fill=BaselineBlue!62,draw=BaselineBlue!45!black,line width=0.08pt] (3,0) rectangle ++(1,-1);
\filldraw[fill=white,draw=black!15,line width=0.08pt] (4,0) rectangle ++(1,-1);
\filldraw[fill=BaselineBlue!88,draw=BaselineBlue!45!black,line width=0.08pt] (5,0) rectangle ++(1,-1);
\filldraw[fill=BaselineBlue!62,draw=BaselineBlue!45!black,line width=0.08pt] (6,0) rectangle ++(1,-1);
\filldraw[fill=InterceptOrange!62,draw=InterceptOrange!55!black,line width=0.08pt] (7,0) rectangle ++(1,-1);
\filldraw[fill=BaselineBlue!62,draw=BaselineBlue!45!black,line width=0.08pt] (8,0) rectangle ++(1,-1);
\filldraw[fill=BaselineBlue!62,draw=BaselineBlue!45!black,line width=0.08pt] (9,0) rectangle ++(1,-1);
\filldraw[fill=BaselineBlue!88,draw=BaselineBlue!45!black,line width=0.08pt] (10,0) rectangle ++(1,-1);
\filldraw[fill=white,draw=black!15,line width=0.08pt] (11,0) rectangle ++(1,-1);
\filldraw[fill=BaselineBlue!62,draw=BaselineBlue!45!black,line width=0.08pt] (12,0) rectangle ++(1,-1);
\filldraw[fill=InterceptOrange!62,draw=InterceptOrange!55!black,line width=0.08pt] (13,0) rectangle ++(1,-1);
\filldraw[fill=white,draw=black!15,line width=0.08pt] (14,0) rectangle ++(1,-1);
\filldraw[fill=BaselineBlue!88,draw=BaselineBlue!45!black,line width=0.08pt] (15,0) rectangle ++(1,-1);
\filldraw[fill=white,draw=black!15,line width=0.08pt] (0,-1) rectangle ++(1,-1);
\filldraw[fill=white,draw=black!15,line width=0.08pt] (1,-1) rectangle ++(1,-1);
\filldraw[fill=white,draw=black!15,line width=0.08pt] (2,-1) rectangle ++(1,-1);
\filldraw[fill=white,draw=black!15,line width=0.08pt] (3,-1) rectangle ++(1,-1);
\filldraw[fill=white,draw=black!15,line width=0.08pt] (4,-1) rectangle ++(1,-1);
\filldraw[fill=white,draw=black!15,line width=0.08pt] (5,-1) rectangle ++(1,-1);
\filldraw[fill=white,draw=black!15,line width=0.08pt] (6,-1) rectangle ++(1,-1);
\filldraw[fill=white,draw=black!15,line width=0.08pt] (7,-1) rectangle ++(1,-1);
\filldraw[fill=white,draw=black!15,line width=0.08pt] (8,-1) rectangle ++(1,-1);
\filldraw[fill=white,draw=black!15,line width=0.08pt] (9,-1) rectangle ++(1,-1);
\filldraw[fill=white,draw=black!15,line width=0.08pt] (10,-1) rectangle ++(1,-1);
\filldraw[fill=white,draw=black!15,line width=0.08pt] (11,-1) rectangle ++(1,-1);
\filldraw[fill=white,draw=black!15,line width=0.08pt] (12,-1) rectangle ++(1,-1);
\filldraw[fill=white,draw=black!15,line width=0.08pt] (13,-1) rectangle ++(1,-1);
\filldraw[fill=white,draw=black!15,line width=0.08pt] (14,-1) rectangle ++(1,-1);
\filldraw[fill=white,draw=black!15,line width=0.08pt] (15,-1) rectangle ++(1,-1);
\filldraw[fill=BaselineBlue!62,draw=BaselineBlue!45!black,line width=0.08pt] (0,-2) rectangle ++(1,-1);
\filldraw[fill=white,draw=black!15,line width=0.08pt] (1,-2) rectangle ++(1,-1);
\filldraw[fill=BaselineBlue!44,draw=BaselineBlue!45!black,line width=0.08pt] (2,-2) rectangle ++(1,-1);
\filldraw[fill=BaselineBlue!44,draw=BaselineBlue!45!black,line width=0.08pt] (3,-2) rectangle ++(1,-1);
\filldraw[fill=white,draw=black!15,line width=0.08pt] (4,-2) rectangle ++(1,-1);
\filldraw[fill=BaselineBlue!62,draw=BaselineBlue!45!black,line width=0.08pt] (5,-2) rectangle ++(1,-1);
\filldraw[fill=BaselineBlue!44,draw=BaselineBlue!45!black,line width=0.08pt] (6,-2) rectangle ++(1,-1);
\filldraw[fill=InterceptOrange!44,draw=InterceptOrange!55!black,line width=0.08pt] (7,-2) rectangle ++(1,-1);
\filldraw[fill=BaselineBlue!44,draw=BaselineBlue!45!black,line width=0.08pt] (8,-2) rectangle ++(1,-1);
\filldraw[fill=BaselineBlue!44,draw=BaselineBlue!45!black,line width=0.08pt] (9,-2) rectangle ++(1,-1);
\filldraw[fill=BaselineBlue!62,draw=BaselineBlue!45!black,line width=0.08pt] (10,-2) rectangle ++(1,-1);
\filldraw[fill=white,draw=black!15,line width=0.08pt] (11,-2) rectangle ++(1,-1);
\filldraw[fill=BaselineBlue!44,draw=BaselineBlue!45!black,line width=0.08pt] (12,-2) rectangle ++(1,-1);
\filldraw[fill=InterceptOrange!44,draw=InterceptOrange!55!black,line width=0.08pt] (13,-2) rectangle ++(1,-1);
\filldraw[fill=white,draw=black!15,line width=0.08pt] (14,-2) rectangle ++(1,-1);
\filldraw[fill=BaselineBlue!62,draw=BaselineBlue!45!black,line width=0.08pt] (15,-2) rectangle ++(1,-1);
\filldraw[fill=BaselineBlue!62,draw=BaselineBlue!45!black,line width=0.08pt] (0,-3) rectangle ++(1,-1);
\filldraw[fill=white,draw=black!15,line width=0.08pt] (1,-3) rectangle ++(1,-1);
\filldraw[fill=BaselineBlue!44,draw=BaselineBlue!45!black,line width=0.08pt] (2,-3) rectangle ++(1,-1);
\filldraw[fill=BaselineBlue!44,draw=BaselineBlue!45!black,line width=0.08pt] (3,-3) rectangle ++(1,-1);
\filldraw[fill=white,draw=black!15,line width=0.08pt] (4,-3) rectangle ++(1,-1);
\filldraw[fill=BaselineBlue!62,draw=BaselineBlue!45!black,line width=0.08pt] (5,-3) rectangle ++(1,-1);
\filldraw[fill=BaselineBlue!44,draw=BaselineBlue!45!black,line width=0.08pt] (6,-3) rectangle ++(1,-1);
\filldraw[fill=InterceptOrange!44,draw=InterceptOrange!55!black,line width=0.08pt] (7,-3) rectangle ++(1,-1);
\filldraw[fill=BaselineBlue!44,draw=BaselineBlue!45!black,line width=0.08pt] (8,-3) rectangle ++(1,-1);
\filldraw[fill=BaselineBlue!44,draw=BaselineBlue!45!black,line width=0.08pt] (9,-3) rectangle ++(1,-1);
\filldraw[fill=BaselineBlue!62,draw=BaselineBlue!45!black,line width=0.08pt] (10,-3) rectangle ++(1,-1);
\filldraw[fill=white,draw=black!15,line width=0.08pt] (11,-3) rectangle ++(1,-1);
\filldraw[fill=BaselineBlue!44,draw=BaselineBlue!45!black,line width=0.08pt] (12,-3) rectangle ++(1,-1);
\filldraw[fill=InterceptOrange!44,draw=InterceptOrange!55!black,line width=0.08pt] (13,-3) rectangle ++(1,-1);
\filldraw[fill=white,draw=black!15,line width=0.08pt] (14,-3) rectangle ++(1,-1);
\filldraw[fill=BaselineBlue!62,draw=BaselineBlue!45!black,line width=0.08pt] (15,-3) rectangle ++(1,-1);
\filldraw[fill=white,draw=black!15,line width=0.08pt] (0,-4) rectangle ++(1,-1);
\filldraw[fill=white,draw=black!15,line width=0.08pt] (1,-4) rectangle ++(1,-1);
\filldraw[fill=white,draw=black!15,line width=0.08pt] (2,-4) rectangle ++(1,-1);
\filldraw[fill=white,draw=black!15,line width=0.08pt] (3,-4) rectangle ++(1,-1);
\filldraw[fill=white,draw=black!15,line width=0.08pt] (4,-4) rectangle ++(1,-1);
\filldraw[fill=white,draw=black!15,line width=0.08pt] (5,-4) rectangle ++(1,-1);
\filldraw[fill=white,draw=black!15,line width=0.08pt] (6,-4) rectangle ++(1,-1);
\filldraw[fill=white,draw=black!15,line width=0.08pt] (7,-4) rectangle ++(1,-1);
\filldraw[fill=white,draw=black!15,line width=0.08pt] (8,-4) rectangle ++(1,-1);
\filldraw[fill=white,draw=black!15,line width=0.08pt] (9,-4) rectangle ++(1,-1);
\filldraw[fill=white,draw=black!15,line width=0.08pt] (10,-4) rectangle ++(1,-1);
\filldraw[fill=white,draw=black!15,line width=0.08pt] (11,-4) rectangle ++(1,-1);
\filldraw[fill=white,draw=black!15,line width=0.08pt] (12,-4) rectangle ++(1,-1);
\filldraw[fill=white,draw=black!15,line width=0.08pt] (13,-4) rectangle ++(1,-1);
\filldraw[fill=white,draw=black!15,line width=0.08pt] (14,-4) rectangle ++(1,-1);
\filldraw[fill=white,draw=black!15,line width=0.08pt] (15,-4) rectangle ++(1,-1);
\filldraw[fill=BaselineBlue!88,draw=BaselineBlue!45!black,line width=0.08pt] (0,-5) rectangle ++(1,-1);
\filldraw[fill=white,draw=black!15,line width=0.08pt] (1,-5) rectangle ++(1,-1);
\filldraw[fill=BaselineBlue!62,draw=BaselineBlue!45!black,line width=0.08pt] (2,-5) rectangle ++(1,-1);
\filldraw[fill=BaselineBlue!62,draw=BaselineBlue!45!black,line width=0.08pt] (3,-5) rectangle ++(1,-1);
\filldraw[fill=white,draw=black!15,line width=0.08pt] (4,-5) rectangle ++(1,-1);
\filldraw[fill=BaselineBlue!88,draw=BaselineBlue!45!black,line width=0.08pt] (5,-5) rectangle ++(1,-1);
\filldraw[fill=BaselineBlue!62,draw=BaselineBlue!45!black,line width=0.08pt] (6,-5) rectangle ++(1,-1);
\filldraw[fill=InterceptOrange!62,draw=InterceptOrange!55!black,line width=0.08pt] (7,-5) rectangle ++(1,-1);
\filldraw[fill=BaselineBlue!62,draw=BaselineBlue!45!black,line width=0.08pt] (8,-5) rectangle ++(1,-1);
\filldraw[fill=BaselineBlue!62,draw=BaselineBlue!45!black,line width=0.08pt] (9,-5) rectangle ++(1,-1);
\filldraw[fill=BaselineBlue!88,draw=BaselineBlue!45!black,line width=0.08pt] (10,-5) rectangle ++(1,-1);
\filldraw[fill=white,draw=black!15,line width=0.08pt] (11,-5) rectangle ++(1,-1);
\filldraw[fill=BaselineBlue!62,draw=BaselineBlue!45!black,line width=0.08pt] (12,-5) rectangle ++(1,-1);
\filldraw[fill=InterceptOrange!62,draw=InterceptOrange!55!black,line width=0.08pt] (13,-5) rectangle ++(1,-1);
\filldraw[fill=white,draw=black!15,line width=0.08pt] (14,-5) rectangle ++(1,-1);
\filldraw[fill=BaselineBlue!88,draw=BaselineBlue!45!black,line width=0.08pt] (15,-5) rectangle ++(1,-1);
\filldraw[fill=BaselineBlue!62,draw=BaselineBlue!45!black,line width=0.08pt] (0,-6) rectangle ++(1,-1);
\filldraw[fill=white,draw=black!15,line width=0.08pt] (1,-6) rectangle ++(1,-1);
\filldraw[fill=BaselineBlue!44,draw=BaselineBlue!45!black,line width=0.08pt] (2,-6) rectangle ++(1,-1);
\filldraw[fill=BaselineBlue!44,draw=BaselineBlue!45!black,line width=0.08pt] (3,-6) rectangle ++(1,-1);
\filldraw[fill=white,draw=black!15,line width=0.08pt] (4,-6) rectangle ++(1,-1);
\filldraw[fill=BaselineBlue!62,draw=BaselineBlue!45!black,line width=0.08pt] (5,-6) rectangle ++(1,-1);
\filldraw[fill=BaselineBlue!44,draw=BaselineBlue!45!black,line width=0.08pt] (6,-6) rectangle ++(1,-1);
\filldraw[fill=InterceptOrange!44,draw=InterceptOrange!55!black,line width=0.08pt] (7,-6) rectangle ++(1,-1);
\filldraw[fill=BaselineBlue!44,draw=BaselineBlue!45!black,line width=0.08pt] (8,-6) rectangle ++(1,-1);
\filldraw[fill=BaselineBlue!44,draw=BaselineBlue!45!black,line width=0.08pt] (9,-6) rectangle ++(1,-1);
\filldraw[fill=BaselineBlue!62,draw=BaselineBlue!45!black,line width=0.08pt] (10,-6) rectangle ++(1,-1);
\filldraw[fill=white,draw=black!15,line width=0.08pt] (11,-6) rectangle ++(1,-1);
\filldraw[fill=BaselineBlue!44,draw=BaselineBlue!45!black,line width=0.08pt] (12,-6) rectangle ++(1,-1);
\filldraw[fill=InterceptOrange!44,draw=InterceptOrange!55!black,line width=0.08pt] (13,-6) rectangle ++(1,-1);
\filldraw[fill=white,draw=black!15,line width=0.08pt] (14,-6) rectangle ++(1,-1);
\filldraw[fill=BaselineBlue!62,draw=BaselineBlue!45!black,line width=0.08pt] (15,-6) rectangle ++(1,-1);
\filldraw[fill=InterceptOrange!62,draw=InterceptOrange!55!black,line width=0.08pt] (0,-7) rectangle ++(1,-1);
\filldraw[fill=white,draw=black!15,line width=0.08pt] (1,-7) rectangle ++(1,-1);
\filldraw[fill=InterceptOrange!44,draw=InterceptOrange!55!black,line width=0.08pt] (2,-7) rectangle ++(1,-1);
\filldraw[fill=InterceptOrange!44,draw=InterceptOrange!55!black,line width=0.08pt] (3,-7) rectangle ++(1,-1);
\filldraw[fill=white,draw=black!15,line width=0.08pt] (4,-7) rectangle ++(1,-1);
\filldraw[fill=InterceptOrange!62,draw=InterceptOrange!55!black,line width=0.08pt] (5,-7) rectangle ++(1,-1);
\filldraw[fill=InterceptOrange!44,draw=InterceptOrange!55!black,line width=0.08pt] (6,-7) rectangle ++(1,-1);
\filldraw[fill=BaselineBlue!44,draw=BaselineBlue!45!black,line width=0.08pt] (7,-7) rectangle ++(1,-1);
\filldraw[fill=InterceptOrange!44,draw=InterceptOrange!55!black,line width=0.08pt] (8,-7) rectangle ++(1,-1);
\filldraw[fill=InterceptOrange!44,draw=InterceptOrange!55!black,line width=0.08pt] (9,-7) rectangle ++(1,-1);
\filldraw[fill=InterceptOrange!62,draw=InterceptOrange!55!black,line width=0.08pt] (10,-7) rectangle ++(1,-1);
\filldraw[fill=white,draw=black!15,line width=0.08pt] (11,-7) rectangle ++(1,-1);
\filldraw[fill=InterceptOrange!44,draw=InterceptOrange!55!black,line width=0.08pt] (12,-7) rectangle ++(1,-1);
\filldraw[fill=BaselineBlue!44,draw=BaselineBlue!45!black,line width=0.08pt] (13,-7) rectangle ++(1,-1);
\filldraw[fill=white,draw=black!15,line width=0.08pt] (14,-7) rectangle ++(1,-1);
\filldraw[fill=InterceptOrange!62,draw=InterceptOrange!55!black,line width=0.08pt] (15,-7) rectangle ++(1,-1);
\filldraw[fill=BaselineBlue!62,draw=BaselineBlue!45!black,line width=0.08pt] (0,-8) rectangle ++(1,-1);
\filldraw[fill=white,draw=black!15,line width=0.08pt] (1,-8) rectangle ++(1,-1);
\filldraw[fill=BaselineBlue!44,draw=BaselineBlue!45!black,line width=0.08pt] (2,-8) rectangle ++(1,-1);
\filldraw[fill=BaselineBlue!44,draw=BaselineBlue!45!black,line width=0.08pt] (3,-8) rectangle ++(1,-1);
\filldraw[fill=white,draw=black!15,line width=0.08pt] (4,-8) rectangle ++(1,-1);
\filldraw[fill=BaselineBlue!62,draw=BaselineBlue!45!black,line width=0.08pt] (5,-8) rectangle ++(1,-1);
\filldraw[fill=BaselineBlue!44,draw=BaselineBlue!45!black,line width=0.08pt] (6,-8) rectangle ++(1,-1);
\filldraw[fill=InterceptOrange!44,draw=InterceptOrange!55!black,line width=0.08pt] (7,-8) rectangle ++(1,-1);
\filldraw[fill=BaselineBlue!44,draw=BaselineBlue!45!black,line width=0.08pt] (8,-8) rectangle ++(1,-1);
\filldraw[fill=BaselineBlue!44,draw=BaselineBlue!45!black,line width=0.08pt] (9,-8) rectangle ++(1,-1);
\filldraw[fill=BaselineBlue!62,draw=BaselineBlue!45!black,line width=0.08pt] (10,-8) rectangle ++(1,-1);
\filldraw[fill=white,draw=black!15,line width=0.08pt] (11,-8) rectangle ++(1,-1);
\filldraw[fill=BaselineBlue!44,draw=BaselineBlue!45!black,line width=0.08pt] (12,-8) rectangle ++(1,-1);
\filldraw[fill=InterceptOrange!44,draw=InterceptOrange!55!black,line width=0.08pt] (13,-8) rectangle ++(1,-1);
\filldraw[fill=white,draw=black!15,line width=0.08pt] (14,-8) rectangle ++(1,-1);
\filldraw[fill=BaselineBlue!62,draw=BaselineBlue!45!black,line width=0.08pt] (15,-8) rectangle ++(1,-1);
\filldraw[fill=BaselineBlue!62,draw=BaselineBlue!45!black,line width=0.08pt] (0,-9) rectangle ++(1,-1);
\filldraw[fill=white,draw=black!15,line width=0.08pt] (1,-9) rectangle ++(1,-1);
\filldraw[fill=BaselineBlue!44,draw=BaselineBlue!45!black,line width=0.08pt] (2,-9) rectangle ++(1,-1);
\filldraw[fill=BaselineBlue!44,draw=BaselineBlue!45!black,line width=0.08pt] (3,-9) rectangle ++(1,-1);
\filldraw[fill=white,draw=black!15,line width=0.08pt] (4,-9) rectangle ++(1,-1);
\filldraw[fill=BaselineBlue!62,draw=BaselineBlue!45!black,line width=0.08pt] (5,-9) rectangle ++(1,-1);
\filldraw[fill=BaselineBlue!44,draw=BaselineBlue!45!black,line width=0.08pt] (6,-9) rectangle ++(1,-1);
\filldraw[fill=InterceptOrange!44,draw=InterceptOrange!55!black,line width=0.08pt] (7,-9) rectangle ++(1,-1);
\filldraw[fill=BaselineBlue!44,draw=BaselineBlue!45!black,line width=0.08pt] (8,-9) rectangle ++(1,-1);
\filldraw[fill=BaselineBlue!44,draw=BaselineBlue!45!black,line width=0.08pt] (9,-9) rectangle ++(1,-1);
\filldraw[fill=BaselineBlue!62,draw=BaselineBlue!45!black,line width=0.08pt] (10,-9) rectangle ++(1,-1);
\filldraw[fill=white,draw=black!15,line width=0.08pt] (11,-9) rectangle ++(1,-1);
\filldraw[fill=BaselineBlue!44,draw=BaselineBlue!45!black,line width=0.08pt] (12,-9) rectangle ++(1,-1);
\filldraw[fill=InterceptOrange!44,draw=InterceptOrange!55!black,line width=0.08pt] (13,-9) rectangle ++(1,-1);
\filldraw[fill=white,draw=black!15,line width=0.08pt] (14,-9) rectangle ++(1,-1);
\filldraw[fill=BaselineBlue!62,draw=BaselineBlue!45!black,line width=0.08pt] (15,-9) rectangle ++(1,-1);
\filldraw[fill=BaselineBlue!88,draw=BaselineBlue!45!black,line width=0.08pt] (0,-10) rectangle ++(1,-1);
\filldraw[fill=white,draw=black!15,line width=0.08pt] (1,-10) rectangle ++(1,-1);
\filldraw[fill=BaselineBlue!62,draw=BaselineBlue!45!black,line width=0.08pt] (2,-10) rectangle ++(1,-1);
\filldraw[fill=BaselineBlue!62,draw=BaselineBlue!45!black,line width=0.08pt] (3,-10) rectangle ++(1,-1);
\filldraw[fill=white,draw=black!15,line width=0.08pt] (4,-10) rectangle ++(1,-1);
\filldraw[fill=BaselineBlue!88,draw=BaselineBlue!45!black,line width=0.08pt] (5,-10) rectangle ++(1,-1);
\filldraw[fill=BaselineBlue!62,draw=BaselineBlue!45!black,line width=0.08pt] (6,-10) rectangle ++(1,-1);
\filldraw[fill=InterceptOrange!62,draw=InterceptOrange!55!black,line width=0.08pt] (7,-10) rectangle ++(1,-1);
\filldraw[fill=BaselineBlue!62,draw=BaselineBlue!45!black,line width=0.08pt] (8,-10) rectangle ++(1,-1);
\filldraw[fill=BaselineBlue!62,draw=BaselineBlue!45!black,line width=0.08pt] (9,-10) rectangle ++(1,-1);
\filldraw[fill=BaselineBlue!88,draw=BaselineBlue!45!black,line width=0.08pt] (10,-10) rectangle ++(1,-1);
\filldraw[fill=white,draw=black!15,line width=0.08pt] (11,-10) rectangle ++(1,-1);
\filldraw[fill=BaselineBlue!62,draw=BaselineBlue!45!black,line width=0.08pt] (12,-10) rectangle ++(1,-1);
\filldraw[fill=InterceptOrange!62,draw=InterceptOrange!55!black,line width=0.08pt] (13,-10) rectangle ++(1,-1);
\filldraw[fill=white,draw=black!15,line width=0.08pt] (14,-10) rectangle ++(1,-1);
\filldraw[fill=BaselineBlue!88,draw=BaselineBlue!45!black,line width=0.08pt] (15,-10) rectangle ++(1,-1);
\filldraw[fill=white,draw=black!15,line width=0.08pt] (0,-11) rectangle ++(1,-1);
\filldraw[fill=white,draw=black!15,line width=0.08pt] (1,-11) rectangle ++(1,-1);
\filldraw[fill=white,draw=black!15,line width=0.08pt] (2,-11) rectangle ++(1,-1);
\filldraw[fill=white,draw=black!15,line width=0.08pt] (3,-11) rectangle ++(1,-1);
\filldraw[fill=white,draw=black!15,line width=0.08pt] (4,-11) rectangle ++(1,-1);
\filldraw[fill=white,draw=black!15,line width=0.08pt] (5,-11) rectangle ++(1,-1);
\filldraw[fill=white,draw=black!15,line width=0.08pt] (6,-11) rectangle ++(1,-1);
\filldraw[fill=white,draw=black!15,line width=0.08pt] (7,-11) rectangle ++(1,-1);
\filldraw[fill=white,draw=black!15,line width=0.08pt] (8,-11) rectangle ++(1,-1);
\filldraw[fill=white,draw=black!15,line width=0.08pt] (9,-11) rectangle ++(1,-1);
\filldraw[fill=white,draw=black!15,line width=0.08pt] (10,-11) rectangle ++(1,-1);
\filldraw[fill=white,draw=black!15,line width=0.08pt] (11,-11) rectangle ++(1,-1);
\filldraw[fill=white,draw=black!15,line width=0.08pt] (12,-11) rectangle ++(1,-1);
\filldraw[fill=white,draw=black!15,line width=0.08pt] (13,-11) rectangle ++(1,-1);
\filldraw[fill=white,draw=black!15,line width=0.08pt] (14,-11) rectangle ++(1,-1);
\filldraw[fill=white,draw=black!15,line width=0.08pt] (15,-11) rectangle ++(1,-1);
\filldraw[fill=BaselineBlue!62,draw=BaselineBlue!45!black,line width=0.08pt] (0,-12) rectangle ++(1,-1);
\filldraw[fill=white,draw=black!15,line width=0.08pt] (1,-12) rectangle ++(1,-1);
\filldraw[fill=BaselineBlue!44,draw=BaselineBlue!45!black,line width=0.08pt] (2,-12) rectangle ++(1,-1);
\filldraw[fill=BaselineBlue!44,draw=BaselineBlue!45!black,line width=0.08pt] (3,-12) rectangle ++(1,-1);
\filldraw[fill=white,draw=black!15,line width=0.08pt] (4,-12) rectangle ++(1,-1);
\filldraw[fill=BaselineBlue!62,draw=BaselineBlue!45!black,line width=0.08pt] (5,-12) rectangle ++(1,-1);
\filldraw[fill=BaselineBlue!44,draw=BaselineBlue!45!black,line width=0.08pt] (6,-12) rectangle ++(1,-1);
\filldraw[fill=InterceptOrange!44,draw=InterceptOrange!55!black,line width=0.08pt] (7,-12) rectangle ++(1,-1);
\filldraw[fill=BaselineBlue!44,draw=BaselineBlue!45!black,line width=0.08pt] (8,-12) rectangle ++(1,-1);
\filldraw[fill=BaselineBlue!44,draw=BaselineBlue!45!black,line width=0.08pt] (9,-12) rectangle ++(1,-1);
\filldraw[fill=BaselineBlue!62,draw=BaselineBlue!45!black,line width=0.08pt] (10,-12) rectangle ++(1,-1);
\filldraw[fill=white,draw=black!15,line width=0.08pt] (11,-12) rectangle ++(1,-1);
\filldraw[fill=BaselineBlue!44,draw=BaselineBlue!45!black,line width=0.08pt] (12,-12) rectangle ++(1,-1);
\filldraw[fill=InterceptOrange!44,draw=InterceptOrange!55!black,line width=0.08pt] (13,-12) rectangle ++(1,-1);
\filldraw[fill=white,draw=black!15,line width=0.08pt] (14,-12) rectangle ++(1,-1);
\filldraw[fill=BaselineBlue!62,draw=BaselineBlue!45!black,line width=0.08pt] (15,-12) rectangle ++(1,-1);
\filldraw[fill=InterceptOrange!62,draw=InterceptOrange!55!black,line width=0.08pt] (0,-13) rectangle ++(1,-1);
\filldraw[fill=white,draw=black!15,line width=0.08pt] (1,-13) rectangle ++(1,-1);
\filldraw[fill=InterceptOrange!44,draw=InterceptOrange!55!black,line width=0.08pt] (2,-13) rectangle ++(1,-1);
\filldraw[fill=InterceptOrange!44,draw=InterceptOrange!55!black,line width=0.08pt] (3,-13) rectangle ++(1,-1);
\filldraw[fill=white,draw=black!15,line width=0.08pt] (4,-13) rectangle ++(1,-1);
\filldraw[fill=InterceptOrange!62,draw=InterceptOrange!55!black,line width=0.08pt] (5,-13) rectangle ++(1,-1);
\filldraw[fill=InterceptOrange!44,draw=InterceptOrange!55!black,line width=0.08pt] (6,-13) rectangle ++(1,-1);
\filldraw[fill=BaselineBlue!44,draw=BaselineBlue!45!black,line width=0.08pt] (7,-13) rectangle ++(1,-1);
\filldraw[fill=InterceptOrange!44,draw=InterceptOrange!55!black,line width=0.08pt] (8,-13) rectangle ++(1,-1);
\filldraw[fill=InterceptOrange!44,draw=InterceptOrange!55!black,line width=0.08pt] (9,-13) rectangle ++(1,-1);
\filldraw[fill=InterceptOrange!62,draw=InterceptOrange!55!black,line width=0.08pt] (10,-13) rectangle ++(1,-1);
\filldraw[fill=white,draw=black!15,line width=0.08pt] (11,-13) rectangle ++(1,-1);
\filldraw[fill=InterceptOrange!44,draw=InterceptOrange!55!black,line width=0.08pt] (12,-13) rectangle ++(1,-1);
\filldraw[fill=BaselineBlue!44,draw=BaselineBlue!45!black,line width=0.08pt] (13,-13) rectangle ++(1,-1);
\filldraw[fill=white,draw=black!15,line width=0.08pt] (14,-13) rectangle ++(1,-1);
\filldraw[fill=InterceptOrange!62,draw=InterceptOrange!55!black,line width=0.08pt] (15,-13) rectangle ++(1,-1);
\filldraw[fill=white,draw=black!15,line width=0.08pt] (0,-14) rectangle ++(1,-1);
\filldraw[fill=white,draw=black!15,line width=0.08pt] (1,-14) rectangle ++(1,-1);
\filldraw[fill=white,draw=black!15,line width=0.08pt] (2,-14) rectangle ++(1,-1);
\filldraw[fill=white,draw=black!15,line width=0.08pt] (3,-14) rectangle ++(1,-1);
\filldraw[fill=white,draw=black!15,line width=0.08pt] (4,-14) rectangle ++(1,-1);
\filldraw[fill=white,draw=black!15,line width=0.08pt] (5,-14) rectangle ++(1,-1);
\filldraw[fill=white,draw=black!15,line width=0.08pt] (6,-14) rectangle ++(1,-1);
\filldraw[fill=white,draw=black!15,line width=0.08pt] (7,-14) rectangle ++(1,-1);
\filldraw[fill=white,draw=black!15,line width=0.08pt] (8,-14) rectangle ++(1,-1);
\filldraw[fill=white,draw=black!15,line width=0.08pt] (9,-14) rectangle ++(1,-1);
\filldraw[fill=white,draw=black!15,line width=0.08pt] (10,-14) rectangle ++(1,-1);
\filldraw[fill=white,draw=black!15,line width=0.08pt] (11,-14) rectangle ++(1,-1);
\filldraw[fill=white,draw=black!15,line width=0.08pt] (12,-14) rectangle ++(1,-1);
\filldraw[fill=white,draw=black!15,line width=0.08pt] (13,-14) rectangle ++(1,-1);
\filldraw[fill=white,draw=black!15,line width=0.08pt] (14,-14) rectangle ++(1,-1);
\filldraw[fill=white,draw=black!15,line width=0.08pt] (15,-14) rectangle ++(1,-1);
\filldraw[fill=BaselineBlue!88,draw=BaselineBlue!45!black,line width=0.08pt] (0,-15) rectangle ++(1,-1);
\filldraw[fill=white,draw=black!15,line width=0.08pt] (1,-15) rectangle ++(1,-1);
\filldraw[fill=BaselineBlue!62,draw=BaselineBlue!45!black,line width=0.08pt] (2,-15) rectangle ++(1,-1);
\filldraw[fill=BaselineBlue!62,draw=BaselineBlue!45!black,line width=0.08pt] (3,-15) rectangle ++(1,-1);
\filldraw[fill=white,draw=black!15,line width=0.08pt] (4,-15) rectangle ++(1,-1);
\filldraw[fill=BaselineBlue!88,draw=BaselineBlue!45!black,line width=0.08pt] (5,-15) rectangle ++(1,-1);
\filldraw[fill=BaselineBlue!62,draw=BaselineBlue!45!black,line width=0.08pt] (6,-15) rectangle ++(1,-1);
\filldraw[fill=InterceptOrange!62,draw=InterceptOrange!55!black,line width=0.08pt] (7,-15) rectangle ++(1,-1);
\filldraw[fill=BaselineBlue!62,draw=BaselineBlue!45!black,line width=0.08pt] (8,-15) rectangle ++(1,-1);
\filldraw[fill=BaselineBlue!62,draw=BaselineBlue!45!black,line width=0.08pt] (9,-15) rectangle ++(1,-1);
\filldraw[fill=BaselineBlue!88,draw=BaselineBlue!45!black,line width=0.08pt] (10,-15) rectangle ++(1,-1);
\filldraw[fill=white,draw=black!15,line width=0.08pt] (11,-15) rectangle ++(1,-1);
\filldraw[fill=BaselineBlue!62,draw=BaselineBlue!45!black,line width=0.08pt] (12,-15) rectangle ++(1,-1);
\filldraw[fill=InterceptOrange!62,draw=InterceptOrange!55!black,line width=0.08pt] (13,-15) rectangle ++(1,-1);
\filldraw[fill=white,draw=black!15,line width=0.08pt] (14,-15) rectangle ++(1,-1);
\filldraw[fill=BaselineBlue!88,draw=BaselineBlue!45!black,line width=0.08pt] (15,-15) rectangle ++(1,-1);
\node[anchor=north] at (0.5,-16.15) {0};
\node[anchor=east] at (-0.15,-0.5) {0};
\node[anchor=north] at (4.5,-16.15) {4};
\node[anchor=east] at (-0.15,-4.5) {4};
\node[anchor=north] at (8.5,-16.15) {8};
\node[anchor=east] at (-0.15,-8.5) {8};
\node[anchor=north] at (12.5,-16.15) {12};
\node[anchor=east] at (-0.15,-12.5) {12};
\node[anchor=north] at (15.5,-16.15) {15};
\node[anchor=east] at (-0.15,-15.5) {15};
\node[anchor=north] at (8,-17.15) {column};
\node[rotate=90] at (-1.15,-8) {row};
\end{scope}
\begin{scope}[xshift=5.35cm]
\node[font=\small\bfseries] at (7.5,1.3) {Intercept--resend $\rho_I$};
\filldraw[fill=BaselineBlue!44,draw=BaselineBlue!45!black,line width=0.08pt] (0,0) rectangle ++(1,-1);
\filldraw[fill=white,draw=black!15,line width=0.08pt] (1,0) rectangle ++(1,-1);
\filldraw[fill=BaselineBlue!31,draw=BaselineBlue!45!black,line width=0.08pt] (2,0) rectangle ++(1,-1);
\filldraw[fill=BaselineBlue!31,draw=BaselineBlue!45!black,line width=0.08pt] (3,0) rectangle ++(1,-1);
\filldraw[fill=BaselineBlue!44,draw=BaselineBlue!45!black,line width=0.08pt] (4,0) rectangle ++(1,-1);
\filldraw[fill=white,draw=black!15,line width=0.08pt] (5,0) rectangle ++(1,-1);
\filldraw[fill=BaselineBlue!31,draw=BaselineBlue!45!black,line width=0.08pt] (6,0) rectangle ++(1,-1);
\filldraw[fill=BaselineBlue!31,draw=BaselineBlue!45!black,line width=0.08pt] (7,0) rectangle ++(1,-1);
\filldraw[fill=BaselineBlue!62,draw=BaselineBlue!45!black,line width=0.08pt] (8,0) rectangle ++(1,-1);
\filldraw[fill=white,draw=black!15,line width=0.08pt] (9,0) rectangle ++(1,-1);
\filldraw[fill=BaselineBlue!44,draw=BaselineBlue!45!black,line width=0.08pt] (10,0) rectangle ++(1,-1);
\filldraw[fill=BaselineBlue!44,draw=BaselineBlue!45!black,line width=0.08pt] (11,0) rectangle ++(1,-1);
\filldraw[fill=white,draw=black!15,line width=0.08pt] (12,0) rectangle ++(1,-1);
\filldraw[fill=white,draw=black!15,line width=0.08pt] (13,0) rectangle ++(1,-1);
\filldraw[fill=white,draw=black!15,line width=0.08pt] (14,0) rectangle ++(1,-1);
\filldraw[fill=white,draw=black!15,line width=0.08pt] (15,0) rectangle ++(1,-1);
\filldraw[fill=white,draw=black!15,line width=0.08pt] (0,-1) rectangle ++(1,-1);
\filldraw[fill=BaselineBlue!44,draw=BaselineBlue!45!black,line width=0.08pt] (1,-1) rectangle ++(1,-1);
\filldraw[fill=BaselineBlue!31,draw=BaselineBlue!45!black,line width=0.08pt] (2,-1) rectangle ++(1,-1);
\filldraw[fill=InterceptOrange!31,draw=InterceptOrange!55!black,line width=0.08pt] (3,-1) rectangle ++(1,-1);
\filldraw[fill=white,draw=black!15,line width=0.08pt] (4,-1) rectangle ++(1,-1);
\filldraw[fill=BaselineBlue!44,draw=BaselineBlue!45!black,line width=0.08pt] (5,-1) rectangle ++(1,-1);
\filldraw[fill=BaselineBlue!31,draw=BaselineBlue!45!black,line width=0.08pt] (6,-1) rectangle ++(1,-1);
\filldraw[fill=InterceptOrange!31,draw=InterceptOrange!55!black,line width=0.08pt] (7,-1) rectangle ++(1,-1);
\filldraw[fill=white,draw=black!15,line width=0.08pt] (8,-1) rectangle ++(1,-1);
\filldraw[fill=BaselineBlue!62,draw=BaselineBlue!45!black,line width=0.08pt] (9,-1) rectangle ++(1,-1);
\filldraw[fill=BaselineBlue!44,draw=BaselineBlue!45!black,line width=0.08pt] (10,-1) rectangle ++(1,-1);
\filldraw[fill=InterceptOrange!44,draw=InterceptOrange!55!black,line width=0.08pt] (11,-1) rectangle ++(1,-1);
\filldraw[fill=white,draw=black!15,line width=0.08pt] (12,-1) rectangle ++(1,-1);
\filldraw[fill=white,draw=black!15,line width=0.08pt] (13,-1) rectangle ++(1,-1);
\filldraw[fill=white,draw=black!15,line width=0.08pt] (14,-1) rectangle ++(1,-1);
\filldraw[fill=white,draw=black!15,line width=0.08pt] (15,-1) rectangle ++(1,-1);
\filldraw[fill=BaselineBlue!31,draw=BaselineBlue!45!black,line width=0.08pt] (0,-2) rectangle ++(1,-1);
\filldraw[fill=BaselineBlue!31,draw=BaselineBlue!45!black,line width=0.08pt] (1,-2) rectangle ++(1,-1);
\filldraw[fill=BaselineBlue!44,draw=BaselineBlue!45!black,line width=0.08pt] (2,-2) rectangle ++(1,-1);
\filldraw[fill=white,draw=black!15,line width=0.08pt] (3,-2) rectangle ++(1,-1);
\filldraw[fill=BaselineBlue!31,draw=BaselineBlue!45!black,line width=0.08pt] (4,-2) rectangle ++(1,-1);
\filldraw[fill=BaselineBlue!31,draw=BaselineBlue!45!black,line width=0.08pt] (5,-2) rectangle ++(1,-1);
\filldraw[fill=BaselineBlue!44,draw=BaselineBlue!45!black,line width=0.08pt] (6,-2) rectangle ++(1,-1);
\filldraw[fill=white,draw=black!15,line width=0.08pt] (7,-2) rectangle ++(1,-1);
\filldraw[fill=BaselineBlue!44,draw=BaselineBlue!45!black,line width=0.08pt] (8,-2) rectangle ++(1,-1);
\filldraw[fill=BaselineBlue!44,draw=BaselineBlue!45!black,line width=0.08pt] (9,-2) rectangle ++(1,-1);
\filldraw[fill=BaselineBlue!62,draw=BaselineBlue!45!black,line width=0.08pt] (10,-2) rectangle ++(1,-1);
\filldraw[fill=white,draw=black!15,line width=0.08pt] (11,-2) rectangle ++(1,-1);
\filldraw[fill=white,draw=black!15,line width=0.08pt] (12,-2) rectangle ++(1,-1);
\filldraw[fill=white,draw=black!15,line width=0.08pt] (13,-2) rectangle ++(1,-1);
\filldraw[fill=white,draw=black!15,line width=0.08pt] (14,-2) rectangle ++(1,-1);
\filldraw[fill=white,draw=black!15,line width=0.08pt] (15,-2) rectangle ++(1,-1);
\filldraw[fill=BaselineBlue!31,draw=BaselineBlue!45!black,line width=0.08pt] (0,-3) rectangle ++(1,-1);
\filldraw[fill=InterceptOrange!31,draw=InterceptOrange!55!black,line width=0.08pt] (1,-3) rectangle ++(1,-1);
\filldraw[fill=white,draw=black!15,line width=0.08pt] (2,-3) rectangle ++(1,-1);
\filldraw[fill=BaselineBlue!44,draw=BaselineBlue!45!black,line width=0.08pt] (3,-3) rectangle ++(1,-1);
\filldraw[fill=BaselineBlue!31,draw=BaselineBlue!45!black,line width=0.08pt] (4,-3) rectangle ++(1,-1);
\filldraw[fill=InterceptOrange!31,draw=InterceptOrange!55!black,line width=0.08pt] (5,-3) rectangle ++(1,-1);
\filldraw[fill=white,draw=black!15,line width=0.08pt] (6,-3) rectangle ++(1,-1);
\filldraw[fill=BaselineBlue!44,draw=BaselineBlue!45!black,line width=0.08pt] (7,-3) rectangle ++(1,-1);
\filldraw[fill=BaselineBlue!44,draw=BaselineBlue!45!black,line width=0.08pt] (8,-3) rectangle ++(1,-1);
\filldraw[fill=InterceptOrange!44,draw=InterceptOrange!55!black,line width=0.08pt] (9,-3) rectangle ++(1,-1);
\filldraw[fill=white,draw=black!15,line width=0.08pt] (10,-3) rectangle ++(1,-1);
\filldraw[fill=BaselineBlue!62,draw=BaselineBlue!45!black,line width=0.08pt] (11,-3) rectangle ++(1,-1);
\filldraw[fill=white,draw=black!15,line width=0.08pt] (12,-3) rectangle ++(1,-1);
\filldraw[fill=white,draw=black!15,line width=0.08pt] (13,-3) rectangle ++(1,-1);
\filldraw[fill=white,draw=black!15,line width=0.08pt] (14,-3) rectangle ++(1,-1);
\filldraw[fill=white,draw=black!15,line width=0.08pt] (15,-3) rectangle ++(1,-1);
\filldraw[fill=BaselineBlue!44,draw=BaselineBlue!45!black,line width=0.08pt] (0,-4) rectangle ++(1,-1);
\filldraw[fill=white,draw=black!15,line width=0.08pt] (1,-4) rectangle ++(1,-1);
\filldraw[fill=BaselineBlue!31,draw=BaselineBlue!45!black,line width=0.08pt] (2,-4) rectangle ++(1,-1);
\filldraw[fill=BaselineBlue!31,draw=BaselineBlue!45!black,line width=0.08pt] (3,-4) rectangle ++(1,-1);
\filldraw[fill=BaselineBlue!44,draw=BaselineBlue!45!black,line width=0.08pt] (4,-4) rectangle ++(1,-1);
\filldraw[fill=white,draw=black!15,line width=0.08pt] (5,-4) rectangle ++(1,-1);
\filldraw[fill=BaselineBlue!31,draw=BaselineBlue!45!black,line width=0.08pt] (6,-4) rectangle ++(1,-1);
\filldraw[fill=BaselineBlue!31,draw=BaselineBlue!45!black,line width=0.08pt] (7,-4) rectangle ++(1,-1);
\filldraw[fill=BaselineBlue!62,draw=BaselineBlue!45!black,line width=0.08pt] (8,-4) rectangle ++(1,-1);
\filldraw[fill=white,draw=black!15,line width=0.08pt] (9,-4) rectangle ++(1,-1);
\filldraw[fill=BaselineBlue!44,draw=BaselineBlue!45!black,line width=0.08pt] (10,-4) rectangle ++(1,-1);
\filldraw[fill=BaselineBlue!44,draw=BaselineBlue!45!black,line width=0.08pt] (11,-4) rectangle ++(1,-1);
\filldraw[fill=white,draw=black!15,line width=0.08pt] (12,-4) rectangle ++(1,-1);
\filldraw[fill=white,draw=black!15,line width=0.08pt] (13,-4) rectangle ++(1,-1);
\filldraw[fill=white,draw=black!15,line width=0.08pt] (14,-4) rectangle ++(1,-1);
\filldraw[fill=white,draw=black!15,line width=0.08pt] (15,-4) rectangle ++(1,-1);
\filldraw[fill=white,draw=black!15,line width=0.08pt] (0,-5) rectangle ++(1,-1);
\filldraw[fill=BaselineBlue!44,draw=BaselineBlue!45!black,line width=0.08pt] (1,-5) rectangle ++(1,-1);
\filldraw[fill=BaselineBlue!31,draw=BaselineBlue!45!black,line width=0.08pt] (2,-5) rectangle ++(1,-1);
\filldraw[fill=InterceptOrange!31,draw=InterceptOrange!55!black,line width=0.08pt] (3,-5) rectangle ++(1,-1);
\filldraw[fill=white,draw=black!15,line width=0.08pt] (4,-5) rectangle ++(1,-1);
\filldraw[fill=BaselineBlue!44,draw=BaselineBlue!45!black,line width=0.08pt] (5,-5) rectangle ++(1,-1);
\filldraw[fill=BaselineBlue!31,draw=BaselineBlue!45!black,line width=0.08pt] (6,-5) rectangle ++(1,-1);
\filldraw[fill=InterceptOrange!31,draw=InterceptOrange!55!black,line width=0.08pt] (7,-5) rectangle ++(1,-1);
\filldraw[fill=white,draw=black!15,line width=0.08pt] (8,-5) rectangle ++(1,-1);
\filldraw[fill=BaselineBlue!62,draw=BaselineBlue!45!black,line width=0.08pt] (9,-5) rectangle ++(1,-1);
\filldraw[fill=BaselineBlue!44,draw=BaselineBlue!45!black,line width=0.08pt] (10,-5) rectangle ++(1,-1);
\filldraw[fill=InterceptOrange!44,draw=InterceptOrange!55!black,line width=0.08pt] (11,-5) rectangle ++(1,-1);
\filldraw[fill=white,draw=black!15,line width=0.08pt] (12,-5) rectangle ++(1,-1);
\filldraw[fill=white,draw=black!15,line width=0.08pt] (13,-5) rectangle ++(1,-1);
\filldraw[fill=white,draw=black!15,line width=0.08pt] (14,-5) rectangle ++(1,-1);
\filldraw[fill=white,draw=black!15,line width=0.08pt] (15,-5) rectangle ++(1,-1);
\filldraw[fill=BaselineBlue!31,draw=BaselineBlue!45!black,line width=0.08pt] (0,-6) rectangle ++(1,-1);
\filldraw[fill=BaselineBlue!31,draw=BaselineBlue!45!black,line width=0.08pt] (1,-6) rectangle ++(1,-1);
\filldraw[fill=BaselineBlue!44,draw=BaselineBlue!45!black,line width=0.08pt] (2,-6) rectangle ++(1,-1);
\filldraw[fill=white,draw=black!15,line width=0.08pt] (3,-6) rectangle ++(1,-1);
\filldraw[fill=BaselineBlue!31,draw=BaselineBlue!45!black,line width=0.08pt] (4,-6) rectangle ++(1,-1);
\filldraw[fill=BaselineBlue!31,draw=BaselineBlue!45!black,line width=0.08pt] (5,-6) rectangle ++(1,-1);
\filldraw[fill=BaselineBlue!44,draw=BaselineBlue!45!black,line width=0.08pt] (6,-6) rectangle ++(1,-1);
\filldraw[fill=white,draw=black!15,line width=0.08pt] (7,-6) rectangle ++(1,-1);
\filldraw[fill=BaselineBlue!44,draw=BaselineBlue!45!black,line width=0.08pt] (8,-6) rectangle ++(1,-1);
\filldraw[fill=BaselineBlue!44,draw=BaselineBlue!45!black,line width=0.08pt] (9,-6) rectangle ++(1,-1);
\filldraw[fill=BaselineBlue!62,draw=BaselineBlue!45!black,line width=0.08pt] (10,-6) rectangle ++(1,-1);
\filldraw[fill=white,draw=black!15,line width=0.08pt] (11,-6) rectangle ++(1,-1);
\filldraw[fill=white,draw=black!15,line width=0.08pt] (12,-6) rectangle ++(1,-1);
\filldraw[fill=white,draw=black!15,line width=0.08pt] (13,-6) rectangle ++(1,-1);
\filldraw[fill=white,draw=black!15,line width=0.08pt] (14,-6) rectangle ++(1,-1);
\filldraw[fill=white,draw=black!15,line width=0.08pt] (15,-6) rectangle ++(1,-1);
\filldraw[fill=BaselineBlue!31,draw=BaselineBlue!45!black,line width=0.08pt] (0,-7) rectangle ++(1,-1);
\filldraw[fill=InterceptOrange!31,draw=InterceptOrange!55!black,line width=0.08pt] (1,-7) rectangle ++(1,-1);
\filldraw[fill=white,draw=black!15,line width=0.08pt] (2,-7) rectangle ++(1,-1);
\filldraw[fill=BaselineBlue!44,draw=BaselineBlue!45!black,line width=0.08pt] (3,-7) rectangle ++(1,-1);
\filldraw[fill=BaselineBlue!31,draw=BaselineBlue!45!black,line width=0.08pt] (4,-7) rectangle ++(1,-1);
\filldraw[fill=InterceptOrange!31,draw=InterceptOrange!55!black,line width=0.08pt] (5,-7) rectangle ++(1,-1);
\filldraw[fill=white,draw=black!15,line width=0.08pt] (6,-7) rectangle ++(1,-1);
\filldraw[fill=BaselineBlue!44,draw=BaselineBlue!45!black,line width=0.08pt] (7,-7) rectangle ++(1,-1);
\filldraw[fill=BaselineBlue!44,draw=BaselineBlue!45!black,line width=0.08pt] (8,-7) rectangle ++(1,-1);
\filldraw[fill=InterceptOrange!44,draw=InterceptOrange!55!black,line width=0.08pt] (9,-7) rectangle ++(1,-1);
\filldraw[fill=white,draw=black!15,line width=0.08pt] (10,-7) rectangle ++(1,-1);
\filldraw[fill=BaselineBlue!62,draw=BaselineBlue!45!black,line width=0.08pt] (11,-7) rectangle ++(1,-1);
\filldraw[fill=white,draw=black!15,line width=0.08pt] (12,-7) rectangle ++(1,-1);
\filldraw[fill=white,draw=black!15,line width=0.08pt] (13,-7) rectangle ++(1,-1);
\filldraw[fill=white,draw=black!15,line width=0.08pt] (14,-7) rectangle ++(1,-1);
\filldraw[fill=white,draw=black!15,line width=0.08pt] (15,-7) rectangle ++(1,-1);
\filldraw[fill=BaselineBlue!62,draw=BaselineBlue!45!black,line width=0.08pt] (0,-8) rectangle ++(1,-1);
\filldraw[fill=white,draw=black!15,line width=0.08pt] (1,-8) rectangle ++(1,-1);
\filldraw[fill=BaselineBlue!44,draw=BaselineBlue!45!black,line width=0.08pt] (2,-8) rectangle ++(1,-1);
\filldraw[fill=BaselineBlue!44,draw=BaselineBlue!45!black,line width=0.08pt] (3,-8) rectangle ++(1,-1);
\filldraw[fill=BaselineBlue!62,draw=BaselineBlue!45!black,line width=0.08pt] (4,-8) rectangle ++(1,-1);
\filldraw[fill=white,draw=black!15,line width=0.08pt] (5,-8) rectangle ++(1,-1);
\filldraw[fill=BaselineBlue!44,draw=BaselineBlue!45!black,line width=0.08pt] (6,-8) rectangle ++(1,-1);
\filldraw[fill=BaselineBlue!44,draw=BaselineBlue!45!black,line width=0.08pt] (7,-8) rectangle ++(1,-1);
\filldraw[fill=BaselineBlue!88,draw=BaselineBlue!45!black,line width=0.08pt] (8,-8) rectangle ++(1,-1);
\filldraw[fill=white,draw=black!15,line width=0.08pt] (9,-8) rectangle ++(1,-1);
\filldraw[fill=BaselineBlue!62,draw=BaselineBlue!45!black,line width=0.08pt] (10,-8) rectangle ++(1,-1);
\filldraw[fill=BaselineBlue!62,draw=BaselineBlue!45!black,line width=0.08pt] (11,-8) rectangle ++(1,-1);
\filldraw[fill=white,draw=black!15,line width=0.08pt] (12,-8) rectangle ++(1,-1);
\filldraw[fill=white,draw=black!15,line width=0.08pt] (13,-8) rectangle ++(1,-1);
\filldraw[fill=white,draw=black!15,line width=0.08pt] (14,-8) rectangle ++(1,-1);
\filldraw[fill=white,draw=black!15,line width=0.08pt] (15,-8) rectangle ++(1,-1);
\filldraw[fill=white,draw=black!15,line width=0.08pt] (0,-9) rectangle ++(1,-1);
\filldraw[fill=BaselineBlue!62,draw=BaselineBlue!45!black,line width=0.08pt] (1,-9) rectangle ++(1,-1);
\filldraw[fill=BaselineBlue!44,draw=BaselineBlue!45!black,line width=0.08pt] (2,-9) rectangle ++(1,-1);
\filldraw[fill=InterceptOrange!44,draw=InterceptOrange!55!black,line width=0.08pt] (3,-9) rectangle ++(1,-1);
\filldraw[fill=white,draw=black!15,line width=0.08pt] (4,-9) rectangle ++(1,-1);
\filldraw[fill=BaselineBlue!62,draw=BaselineBlue!45!black,line width=0.08pt] (5,-9) rectangle ++(1,-1);
\filldraw[fill=BaselineBlue!44,draw=BaselineBlue!45!black,line width=0.08pt] (6,-9) rectangle ++(1,-1);
\filldraw[fill=InterceptOrange!44,draw=InterceptOrange!55!black,line width=0.08pt] (7,-9) rectangle ++(1,-1);
\filldraw[fill=white,draw=black!15,line width=0.08pt] (8,-9) rectangle ++(1,-1);
\filldraw[fill=BaselineBlue!88,draw=BaselineBlue!45!black,line width=0.08pt] (9,-9) rectangle ++(1,-1);
\filldraw[fill=BaselineBlue!62,draw=BaselineBlue!45!black,line width=0.08pt] (10,-9) rectangle ++(1,-1);
\filldraw[fill=InterceptOrange!62,draw=InterceptOrange!55!black,line width=0.08pt] (11,-9) rectangle ++(1,-1);
\filldraw[fill=white,draw=black!15,line width=0.08pt] (12,-9) rectangle ++(1,-1);
\filldraw[fill=white,draw=black!15,line width=0.08pt] (13,-9) rectangle ++(1,-1);
\filldraw[fill=white,draw=black!15,line width=0.08pt] (14,-9) rectangle ++(1,-1);
\filldraw[fill=white,draw=black!15,line width=0.08pt] (15,-9) rectangle ++(1,-1);
\filldraw[fill=BaselineBlue!44,draw=BaselineBlue!45!black,line width=0.08pt] (0,-10) rectangle ++(1,-1);
\filldraw[fill=BaselineBlue!44,draw=BaselineBlue!45!black,line width=0.08pt] (1,-10) rectangle ++(1,-1);
\filldraw[fill=BaselineBlue!62,draw=BaselineBlue!45!black,line width=0.08pt] (2,-10) rectangle ++(1,-1);
\filldraw[fill=white,draw=black!15,line width=0.08pt] (3,-10) rectangle ++(1,-1);
\filldraw[fill=BaselineBlue!44,draw=BaselineBlue!45!black,line width=0.08pt] (4,-10) rectangle ++(1,-1);
\filldraw[fill=BaselineBlue!44,draw=BaselineBlue!45!black,line width=0.08pt] (5,-10) rectangle ++(1,-1);
\filldraw[fill=BaselineBlue!62,draw=BaselineBlue!45!black,line width=0.08pt] (6,-10) rectangle ++(1,-1);
\filldraw[fill=white,draw=black!15,line width=0.08pt] (7,-10) rectangle ++(1,-1);
\filldraw[fill=BaselineBlue!62,draw=BaselineBlue!45!black,line width=0.08pt] (8,-10) rectangle ++(1,-1);
\filldraw[fill=BaselineBlue!62,draw=BaselineBlue!45!black,line width=0.08pt] (9,-10) rectangle ++(1,-1);
\filldraw[fill=BaselineBlue!88,draw=BaselineBlue!45!black,line width=0.08pt] (10,-10) rectangle ++(1,-1);
\filldraw[fill=white,draw=black!15,line width=0.08pt] (11,-10) rectangle ++(1,-1);
\filldraw[fill=white,draw=black!15,line width=0.08pt] (12,-10) rectangle ++(1,-1);
\filldraw[fill=white,draw=black!15,line width=0.08pt] (13,-10) rectangle ++(1,-1);
\filldraw[fill=white,draw=black!15,line width=0.08pt] (14,-10) rectangle ++(1,-1);
\filldraw[fill=white,draw=black!15,line width=0.08pt] (15,-10) rectangle ++(1,-1);
\filldraw[fill=BaselineBlue!44,draw=BaselineBlue!45!black,line width=0.08pt] (0,-11) rectangle ++(1,-1);
\filldraw[fill=InterceptOrange!44,draw=InterceptOrange!55!black,line width=0.08pt] (1,-11) rectangle ++(1,-1);
\filldraw[fill=white,draw=black!15,line width=0.08pt] (2,-11) rectangle ++(1,-1);
\filldraw[fill=BaselineBlue!62,draw=BaselineBlue!45!black,line width=0.08pt] (3,-11) rectangle ++(1,-1);
\filldraw[fill=BaselineBlue!44,draw=BaselineBlue!45!black,line width=0.08pt] (4,-11) rectangle ++(1,-1);
\filldraw[fill=InterceptOrange!44,draw=InterceptOrange!55!black,line width=0.08pt] (5,-11) rectangle ++(1,-1);
\filldraw[fill=white,draw=black!15,line width=0.08pt] (6,-11) rectangle ++(1,-1);
\filldraw[fill=BaselineBlue!62,draw=BaselineBlue!45!black,line width=0.08pt] (7,-11) rectangle ++(1,-1);
\filldraw[fill=BaselineBlue!62,draw=BaselineBlue!45!black,line width=0.08pt] (8,-11) rectangle ++(1,-1);
\filldraw[fill=InterceptOrange!62,draw=InterceptOrange!55!black,line width=0.08pt] (9,-11) rectangle ++(1,-1);
\filldraw[fill=white,draw=black!15,line width=0.08pt] (10,-11) rectangle ++(1,-1);
\filldraw[fill=BaselineBlue!88,draw=BaselineBlue!45!black,line width=0.08pt] (11,-11) rectangle ++(1,-1);
\filldraw[fill=white,draw=black!15,line width=0.08pt] (12,-11) rectangle ++(1,-1);
\filldraw[fill=white,draw=black!15,line width=0.08pt] (13,-11) rectangle ++(1,-1);
\filldraw[fill=white,draw=black!15,line width=0.08pt] (14,-11) rectangle ++(1,-1);
\filldraw[fill=white,draw=black!15,line width=0.08pt] (15,-11) rectangle ++(1,-1);
\filldraw[fill=white,draw=black!15,line width=0.08pt] (0,-12) rectangle ++(1,-1);
\filldraw[fill=white,draw=black!15,line width=0.08pt] (1,-12) rectangle ++(1,-1);
\filldraw[fill=white,draw=black!15,line width=0.08pt] (2,-12) rectangle ++(1,-1);
\filldraw[fill=white,draw=black!15,line width=0.08pt] (3,-12) rectangle ++(1,-1);
\filldraw[fill=white,draw=black!15,line width=0.08pt] (4,-12) rectangle ++(1,-1);
\filldraw[fill=white,draw=black!15,line width=0.08pt] (5,-12) rectangle ++(1,-1);
\filldraw[fill=white,draw=black!15,line width=0.08pt] (6,-12) rectangle ++(1,-1);
\filldraw[fill=white,draw=black!15,line width=0.08pt] (7,-12) rectangle ++(1,-1);
\filldraw[fill=white,draw=black!15,line width=0.08pt] (8,-12) rectangle ++(1,-1);
\filldraw[fill=white,draw=black!15,line width=0.08pt] (9,-12) rectangle ++(1,-1);
\filldraw[fill=white,draw=black!15,line width=0.08pt] (10,-12) rectangle ++(1,-1);
\filldraw[fill=white,draw=black!15,line width=0.08pt] (11,-12) rectangle ++(1,-1);
\filldraw[fill=white,draw=black!15,line width=0.08pt] (12,-12) rectangle ++(1,-1);
\filldraw[fill=white,draw=black!15,line width=0.08pt] (13,-12) rectangle ++(1,-1);
\filldraw[fill=white,draw=black!15,line width=0.08pt] (14,-12) rectangle ++(1,-1);
\filldraw[fill=white,draw=black!15,line width=0.08pt] (15,-12) rectangle ++(1,-1);
\filldraw[fill=white,draw=black!15,line width=0.08pt] (0,-13) rectangle ++(1,-1);
\filldraw[fill=white,draw=black!15,line width=0.08pt] (1,-13) rectangle ++(1,-1);
\filldraw[fill=white,draw=black!15,line width=0.08pt] (2,-13) rectangle ++(1,-1);
\filldraw[fill=white,draw=black!15,line width=0.08pt] (3,-13) rectangle ++(1,-1);
\filldraw[fill=white,draw=black!15,line width=0.08pt] (4,-13) rectangle ++(1,-1);
\filldraw[fill=white,draw=black!15,line width=0.08pt] (5,-13) rectangle ++(1,-1);
\filldraw[fill=white,draw=black!15,line width=0.08pt] (6,-13) rectangle ++(1,-1);
\filldraw[fill=white,draw=black!15,line width=0.08pt] (7,-13) rectangle ++(1,-1);
\filldraw[fill=white,draw=black!15,line width=0.08pt] (8,-13) rectangle ++(1,-1);
\filldraw[fill=white,draw=black!15,line width=0.08pt] (9,-13) rectangle ++(1,-1);
\filldraw[fill=white,draw=black!15,line width=0.08pt] (10,-13) rectangle ++(1,-1);
\filldraw[fill=white,draw=black!15,line width=0.08pt] (11,-13) rectangle ++(1,-1);
\filldraw[fill=white,draw=black!15,line width=0.08pt] (12,-13) rectangle ++(1,-1);
\filldraw[fill=white,draw=black!15,line width=0.08pt] (13,-13) rectangle ++(1,-1);
\filldraw[fill=white,draw=black!15,line width=0.08pt] (14,-13) rectangle ++(1,-1);
\filldraw[fill=white,draw=black!15,line width=0.08pt] (15,-13) rectangle ++(1,-1);
\filldraw[fill=white,draw=black!15,line width=0.08pt] (0,-14) rectangle ++(1,-1);
\filldraw[fill=white,draw=black!15,line width=0.08pt] (1,-14) rectangle ++(1,-1);
\filldraw[fill=white,draw=black!15,line width=0.08pt] (2,-14) rectangle ++(1,-1);
\filldraw[fill=white,draw=black!15,line width=0.08pt] (3,-14) rectangle ++(1,-1);
\filldraw[fill=white,draw=black!15,line width=0.08pt] (4,-14) rectangle ++(1,-1);
\filldraw[fill=white,draw=black!15,line width=0.08pt] (5,-14) rectangle ++(1,-1);
\filldraw[fill=white,draw=black!15,line width=0.08pt] (6,-14) rectangle ++(1,-1);
\filldraw[fill=white,draw=black!15,line width=0.08pt] (7,-14) rectangle ++(1,-1);
\filldraw[fill=white,draw=black!15,line width=0.08pt] (8,-14) rectangle ++(1,-1);
\filldraw[fill=white,draw=black!15,line width=0.08pt] (9,-14) rectangle ++(1,-1);
\filldraw[fill=white,draw=black!15,line width=0.08pt] (10,-14) rectangle ++(1,-1);
\filldraw[fill=white,draw=black!15,line width=0.08pt] (11,-14) rectangle ++(1,-1);
\filldraw[fill=white,draw=black!15,line width=0.08pt] (12,-14) rectangle ++(1,-1);
\filldraw[fill=white,draw=black!15,line width=0.08pt] (13,-14) rectangle ++(1,-1);
\filldraw[fill=white,draw=black!15,line width=0.08pt] (14,-14) rectangle ++(1,-1);
\filldraw[fill=white,draw=black!15,line width=0.08pt] (15,-14) rectangle ++(1,-1);
\filldraw[fill=white,draw=black!15,line width=0.08pt] (0,-15) rectangle ++(1,-1);
\filldraw[fill=white,draw=black!15,line width=0.08pt] (1,-15) rectangle ++(1,-1);
\filldraw[fill=white,draw=black!15,line width=0.08pt] (2,-15) rectangle ++(1,-1);
\filldraw[fill=white,draw=black!15,line width=0.08pt] (3,-15) rectangle ++(1,-1);
\filldraw[fill=white,draw=black!15,line width=0.08pt] (4,-15) rectangle ++(1,-1);
\filldraw[fill=white,draw=black!15,line width=0.08pt] (5,-15) rectangle ++(1,-1);
\filldraw[fill=white,draw=black!15,line width=0.08pt] (6,-15) rectangle ++(1,-1);
\filldraw[fill=white,draw=black!15,line width=0.08pt] (7,-15) rectangle ++(1,-1);
\filldraw[fill=white,draw=black!15,line width=0.08pt] (8,-15) rectangle ++(1,-1);
\filldraw[fill=white,draw=black!15,line width=0.08pt] (9,-15) rectangle ++(1,-1);
\filldraw[fill=white,draw=black!15,line width=0.08pt] (10,-15) rectangle ++(1,-1);
\filldraw[fill=white,draw=black!15,line width=0.08pt] (11,-15) rectangle ++(1,-1);
\filldraw[fill=white,draw=black!15,line width=0.08pt] (12,-15) rectangle ++(1,-1);
\filldraw[fill=white,draw=black!15,line width=0.08pt] (13,-15) rectangle ++(1,-1);
\filldraw[fill=white,draw=black!15,line width=0.08pt] (14,-15) rectangle ++(1,-1);
\filldraw[fill=white,draw=black!15,line width=0.08pt] (15,-15) rectangle ++(1,-1);
\node[anchor=north] at (0.5,-16.15) {0};
\node[anchor=east] at (-0.15,-0.5) {0};
\node[anchor=north] at (4.5,-16.15) {4};
\node[anchor=east] at (-0.15,-4.5) {4};
\node[anchor=north] at (8.5,-16.15) {8};
\node[anchor=east] at (-0.15,-8.5) {8};
\node[anchor=north] at (12.5,-16.15) {12};
\node[anchor=east] at (-0.15,-12.5) {12};
\node[anchor=north] at (15.5,-16.15) {15};
\node[anchor=east] at (-0.15,-15.5) {15};
\node[anchor=north] at (8,-17.15) {column};
\node[rotate=90] at (-1.15,-8) {row};
\end{scope}
\begin{scope}[xshift=10.7cm]
\node[font=\small\bfseries] at (7.5,1.3) {Advanced $\rho_A$};
\filldraw[fill=BaselineBlue!88,draw=BaselineBlue!45!black,line width=0.08pt] (0,0) rectangle ++(1,-1);
\filldraw[fill=white,draw=black!15,line width=0.08pt] (1,0) rectangle ++(1,-1);
\filldraw[fill=BaselineBlue!44,draw=BaselineBlue!45!black,line width=0.08pt] (2,0) rectangle ++(1,-1);
\filldraw[fill=BaselineBlue!44,draw=BaselineBlue!45!black,line width=0.08pt] (3,0) rectangle ++(1,-1);
\filldraw[fill=white,draw=black!15,line width=0.08pt] (4,0) rectangle ++(1,-1);
\filldraw[fill=white,draw=black!15,line width=0.08pt] (5,0) rectangle ++(1,-1);
\filldraw[fill=BaselineBlue!44,draw=BaselineBlue!45!black,line width=0.08pt] (6,0) rectangle ++(1,-1);
\filldraw[fill=InterceptOrange!44,draw=InterceptOrange!55!black,line width=0.08pt] (7,0) rectangle ++(1,-1);
\filldraw[fill=BaselineBlue!62,draw=BaselineBlue!45!black,line width=0.08pt] (8,0) rectangle ++(1,-1);
\filldraw[fill=white,draw=black!15,line width=0.08pt] (9,0) rectangle ++(1,-1);
\filldraw[fill=BaselineBlue!62,draw=BaselineBlue!45!black,line width=0.08pt] (10,0) rectangle ++(1,-1);
\filldraw[fill=white,draw=black!15,line width=0.08pt] (11,0) rectangle ++(1,-1);
\filldraw[fill=BaselineBlue!62,draw=BaselineBlue!45!black,line width=0.08pt] (12,0) rectangle ++(1,-1);
\filldraw[fill=white,draw=black!15,line width=0.08pt] (13,0) rectangle ++(1,-1);
\filldraw[fill=white,draw=black!15,line width=0.08pt] (14,0) rectangle ++(1,-1);
\filldraw[fill=BaselineBlue!62,draw=BaselineBlue!45!black,line width=0.08pt] (15,0) rectangle ++(1,-1);
\filldraw[fill=white,draw=black!15,line width=0.08pt] (0,-1) rectangle ++(1,-1);
\filldraw[fill=white,draw=black!15,line width=0.08pt] (1,-1) rectangle ++(1,-1);
\filldraw[fill=white,draw=black!15,line width=0.08pt] (2,-1) rectangle ++(1,-1);
\filldraw[fill=white,draw=black!15,line width=0.08pt] (3,-1) rectangle ++(1,-1);
\filldraw[fill=white,draw=black!15,line width=0.08pt] (4,-1) rectangle ++(1,-1);
\filldraw[fill=white,draw=black!15,line width=0.08pt] (5,-1) rectangle ++(1,-1);
\filldraw[fill=white,draw=black!15,line width=0.08pt] (6,-1) rectangle ++(1,-1);
\filldraw[fill=white,draw=black!15,line width=0.08pt] (7,-1) rectangle ++(1,-1);
\filldraw[fill=white,draw=black!15,line width=0.08pt] (8,-1) rectangle ++(1,-1);
\filldraw[fill=white,draw=black!15,line width=0.08pt] (9,-1) rectangle ++(1,-1);
\filldraw[fill=white,draw=black!15,line width=0.08pt] (10,-1) rectangle ++(1,-1);
\filldraw[fill=white,draw=black!15,line width=0.08pt] (11,-1) rectangle ++(1,-1);
\filldraw[fill=white,draw=black!15,line width=0.08pt] (12,-1) rectangle ++(1,-1);
\filldraw[fill=white,draw=black!15,line width=0.08pt] (13,-1) rectangle ++(1,-1);
\filldraw[fill=white,draw=black!15,line width=0.08pt] (14,-1) rectangle ++(1,-1);
\filldraw[fill=white,draw=black!15,line width=0.08pt] (15,-1) rectangle ++(1,-1);
\filldraw[fill=BaselineBlue!44,draw=BaselineBlue!45!black,line width=0.08pt] (0,-2) rectangle ++(1,-1);
\filldraw[fill=white,draw=black!15,line width=0.08pt] (1,-2) rectangle ++(1,-1);
\filldraw[fill=BaselineBlue!44,draw=BaselineBlue!45!black,line width=0.08pt] (2,-2) rectangle ++(1,-1);
\filldraw[fill=white,draw=black!15,line width=0.08pt] (3,-2) rectangle ++(1,-1);
\filldraw[fill=white,draw=black!15,line width=0.08pt] (4,-2) rectangle ++(1,-1);
\filldraw[fill=BaselineBlue!44,draw=BaselineBlue!45!black,line width=0.08pt] (5,-2) rectangle ++(1,-1);
\filldraw[fill=BaselineBlue!44,draw=BaselineBlue!45!black,line width=0.08pt] (6,-2) rectangle ++(1,-1);
\filldraw[fill=white,draw=black!15,line width=0.08pt] (7,-2) rectangle ++(1,-1);
\filldraw[fill=BaselineBlue!31,draw=BaselineBlue!45!black,line width=0.08pt] (8,-2) rectangle ++(1,-1);
\filldraw[fill=BaselineBlue!31,draw=BaselineBlue!45!black,line width=0.08pt] (9,-2) rectangle ++(1,-1);
\filldraw[fill=BaselineBlue!62,draw=BaselineBlue!45!black,line width=0.08pt] (10,-2) rectangle ++(1,-1);
\filldraw[fill=white,draw=black!15,line width=0.08pt] (11,-2) rectangle ++(1,-1);
\filldraw[fill=BaselineBlue!31,draw=BaselineBlue!45!black,line width=0.08pt] (12,-2) rectangle ++(1,-1);
\filldraw[fill=InterceptOrange!31,draw=InterceptOrange!55!black,line width=0.08pt] (13,-2) rectangle ++(1,-1);
\filldraw[fill=white,draw=black!15,line width=0.08pt] (14,-2) rectangle ++(1,-1);
\filldraw[fill=white,draw=black!15,line width=0.08pt] (15,-2) rectangle ++(1,-1);
\filldraw[fill=BaselineBlue!44,draw=BaselineBlue!45!black,line width=0.08pt] (0,-3) rectangle ++(1,-1);
\filldraw[fill=white,draw=black!15,line width=0.08pt] (1,-3) rectangle ++(1,-1);
\filldraw[fill=white,draw=black!15,line width=0.08pt] (2,-3) rectangle ++(1,-1);
\filldraw[fill=BaselineBlue!44,draw=BaselineBlue!45!black,line width=0.08pt] (3,-3) rectangle ++(1,-1);
\filldraw[fill=white,draw=black!15,line width=0.08pt] (4,-3) rectangle ++(1,-1);
\filldraw[fill=InterceptOrange!44,draw=InterceptOrange!55!black,line width=0.08pt] (5,-3) rectangle ++(1,-1);
\filldraw[fill=white,draw=black!15,line width=0.08pt] (6,-3) rectangle ++(1,-1);
\filldraw[fill=InterceptOrange!44,draw=InterceptOrange!55!black,line width=0.08pt] (7,-3) rectangle ++(1,-1);
\filldraw[fill=BaselineBlue!31,draw=BaselineBlue!45!black,line width=0.08pt] (8,-3) rectangle ++(1,-1);
\filldraw[fill=InterceptOrange!31,draw=InterceptOrange!55!black,line width=0.08pt] (9,-3) rectangle ++(1,-1);
\filldraw[fill=white,draw=black!15,line width=0.08pt] (10,-3) rectangle ++(1,-1);
\filldraw[fill=white,draw=black!15,line width=0.08pt] (11,-3) rectangle ++(1,-1);
\filldraw[fill=BaselineBlue!31,draw=BaselineBlue!45!black,line width=0.08pt] (12,-3) rectangle ++(1,-1);
\filldraw[fill=BaselineBlue!31,draw=BaselineBlue!45!black,line width=0.08pt] (13,-3) rectangle ++(1,-1);
\filldraw[fill=white,draw=black!15,line width=0.08pt] (14,-3) rectangle ++(1,-1);
\filldraw[fill=BaselineBlue!62,draw=BaselineBlue!45!black,line width=0.08pt] (15,-3) rectangle ++(1,-1);
\filldraw[fill=white,draw=black!15,line width=0.08pt] (0,-4) rectangle ++(1,-1);
\filldraw[fill=white,draw=black!15,line width=0.08pt] (1,-4) rectangle ++(1,-1);
\filldraw[fill=white,draw=black!15,line width=0.08pt] (2,-4) rectangle ++(1,-1);
\filldraw[fill=white,draw=black!15,line width=0.08pt] (3,-4) rectangle ++(1,-1);
\filldraw[fill=white,draw=black!15,line width=0.08pt] (4,-4) rectangle ++(1,-1);
\filldraw[fill=white,draw=black!15,line width=0.08pt] (5,-4) rectangle ++(1,-1);
\filldraw[fill=white,draw=black!15,line width=0.08pt] (6,-4) rectangle ++(1,-1);
\filldraw[fill=white,draw=black!15,line width=0.08pt] (7,-4) rectangle ++(1,-1);
\filldraw[fill=white,draw=black!15,line width=0.08pt] (8,-4) rectangle ++(1,-1);
\filldraw[fill=white,draw=black!15,line width=0.08pt] (9,-4) rectangle ++(1,-1);
\filldraw[fill=white,draw=black!15,line width=0.08pt] (10,-4) rectangle ++(1,-1);
\filldraw[fill=white,draw=black!15,line width=0.08pt] (11,-4) rectangle ++(1,-1);
\filldraw[fill=white,draw=black!15,line width=0.08pt] (12,-4) rectangle ++(1,-1);
\filldraw[fill=white,draw=black!15,line width=0.08pt] (13,-4) rectangle ++(1,-1);
\filldraw[fill=white,draw=black!15,line width=0.08pt] (14,-4) rectangle ++(1,-1);
\filldraw[fill=white,draw=black!15,line width=0.08pt] (15,-4) rectangle ++(1,-1);
\filldraw[fill=white,draw=black!15,line width=0.08pt] (0,-5) rectangle ++(1,-1);
\filldraw[fill=white,draw=black!15,line width=0.08pt] (1,-5) rectangle ++(1,-1);
\filldraw[fill=BaselineBlue!44,draw=BaselineBlue!45!black,line width=0.08pt] (2,-5) rectangle ++(1,-1);
\filldraw[fill=InterceptOrange!44,draw=InterceptOrange!55!black,line width=0.08pt] (3,-5) rectangle ++(1,-1);
\filldraw[fill=white,draw=black!15,line width=0.08pt] (4,-5) rectangle ++(1,-1);
\filldraw[fill=BaselineBlue!88,draw=BaselineBlue!45!black,line width=0.08pt] (5,-5) rectangle ++(1,-1);
\filldraw[fill=BaselineBlue!44,draw=BaselineBlue!45!black,line width=0.08pt] (6,-5) rectangle ++(1,-1);
\filldraw[fill=BaselineBlue!44,draw=BaselineBlue!45!black,line width=0.08pt] (7,-5) rectangle ++(1,-1);
\filldraw[fill=white,draw=black!15,line width=0.08pt] (8,-5) rectangle ++(1,-1);
\filldraw[fill=BaselineBlue!62,draw=BaselineBlue!45!black,line width=0.08pt] (9,-5) rectangle ++(1,-1);
\filldraw[fill=BaselineBlue!62,draw=BaselineBlue!45!black,line width=0.08pt] (10,-5) rectangle ++(1,-1);
\filldraw[fill=white,draw=black!15,line width=0.08pt] (11,-5) rectangle ++(1,-1);
\filldraw[fill=white,draw=black!15,line width=0.08pt] (12,-5) rectangle ++(1,-1);
\filldraw[fill=InterceptOrange!62,draw=InterceptOrange!55!black,line width=0.08pt] (13,-5) rectangle ++(1,-1);
\filldraw[fill=white,draw=black!15,line width=0.08pt] (14,-5) rectangle ++(1,-1);
\filldraw[fill=InterceptOrange!62,draw=InterceptOrange!55!black,line width=0.08pt] (15,-5) rectangle ++(1,-1);
\filldraw[fill=BaselineBlue!44,draw=BaselineBlue!45!black,line width=0.08pt] (0,-6) rectangle ++(1,-1);
\filldraw[fill=white,draw=black!15,line width=0.08pt] (1,-6) rectangle ++(1,-1);
\filldraw[fill=BaselineBlue!44,draw=BaselineBlue!45!black,line width=0.08pt] (2,-6) rectangle ++(1,-1);
\filldraw[fill=white,draw=black!15,line width=0.08pt] (3,-6) rectangle ++(1,-1);
\filldraw[fill=white,draw=black!15,line width=0.08pt] (4,-6) rectangle ++(1,-1);
\filldraw[fill=BaselineBlue!44,draw=BaselineBlue!45!black,line width=0.08pt] (5,-6) rectangle ++(1,-1);
\filldraw[fill=BaselineBlue!44,draw=BaselineBlue!45!black,line width=0.08pt] (6,-6) rectangle ++(1,-1);
\filldraw[fill=white,draw=black!15,line width=0.08pt] (7,-6) rectangle ++(1,-1);
\filldraw[fill=BaselineBlue!31,draw=BaselineBlue!45!black,line width=0.08pt] (8,-6) rectangle ++(1,-1);
\filldraw[fill=BaselineBlue!31,draw=BaselineBlue!45!black,line width=0.08pt] (9,-6) rectangle ++(1,-1);
\filldraw[fill=BaselineBlue!62,draw=BaselineBlue!45!black,line width=0.08pt] (10,-6) rectangle ++(1,-1);
\filldraw[fill=white,draw=black!15,line width=0.08pt] (11,-6) rectangle ++(1,-1);
\filldraw[fill=BaselineBlue!31,draw=BaselineBlue!45!black,line width=0.08pt] (12,-6) rectangle ++(1,-1);
\filldraw[fill=InterceptOrange!31,draw=InterceptOrange!55!black,line width=0.08pt] (13,-6) rectangle ++(1,-1);
\filldraw[fill=white,draw=black!15,line width=0.08pt] (14,-6) rectangle ++(1,-1);
\filldraw[fill=white,draw=black!15,line width=0.08pt] (15,-6) rectangle ++(1,-1);
\filldraw[fill=InterceptOrange!44,draw=InterceptOrange!55!black,line width=0.08pt] (0,-7) rectangle ++(1,-1);
\filldraw[fill=white,draw=black!15,line width=0.08pt] (1,-7) rectangle ++(1,-1);
\filldraw[fill=white,draw=black!15,line width=0.08pt] (2,-7) rectangle ++(1,-1);
\filldraw[fill=InterceptOrange!44,draw=InterceptOrange!55!black,line width=0.08pt] (3,-7) rectangle ++(1,-1);
\filldraw[fill=white,draw=black!15,line width=0.08pt] (4,-7) rectangle ++(1,-1);
\filldraw[fill=BaselineBlue!44,draw=BaselineBlue!45!black,line width=0.08pt] (5,-7) rectangle ++(1,-1);
\filldraw[fill=white,draw=black!15,line width=0.08pt] (6,-7) rectangle ++(1,-1);
\filldraw[fill=BaselineBlue!44,draw=BaselineBlue!45!black,line width=0.08pt] (7,-7) rectangle ++(1,-1);
\filldraw[fill=InterceptOrange!31,draw=InterceptOrange!55!black,line width=0.08pt] (8,-7) rectangle ++(1,-1);
\filldraw[fill=BaselineBlue!31,draw=BaselineBlue!45!black,line width=0.08pt] (9,-7) rectangle ++(1,-1);
\filldraw[fill=white,draw=black!15,line width=0.08pt] (10,-7) rectangle ++(1,-1);
\filldraw[fill=white,draw=black!15,line width=0.08pt] (11,-7) rectangle ++(1,-1);
\filldraw[fill=InterceptOrange!31,draw=InterceptOrange!55!black,line width=0.08pt] (12,-7) rectangle ++(1,-1);
\filldraw[fill=InterceptOrange!31,draw=InterceptOrange!55!black,line width=0.08pt] (13,-7) rectangle ++(1,-1);
\filldraw[fill=white,draw=black!15,line width=0.08pt] (14,-7) rectangle ++(1,-1);
\filldraw[fill=InterceptOrange!62,draw=InterceptOrange!55!black,line width=0.08pt] (15,-7) rectangle ++(1,-1);
\filldraw[fill=BaselineBlue!62,draw=BaselineBlue!45!black,line width=0.08pt] (0,-8) rectangle ++(1,-1);
\filldraw[fill=white,draw=black!15,line width=0.08pt] (1,-8) rectangle ++(1,-1);
\filldraw[fill=BaselineBlue!31,draw=BaselineBlue!45!black,line width=0.08pt] (2,-8) rectangle ++(1,-1);
\filldraw[fill=BaselineBlue!31,draw=BaselineBlue!45!black,line width=0.08pt] (3,-8) rectangle ++(1,-1);
\filldraw[fill=white,draw=black!15,line width=0.08pt] (4,-8) rectangle ++(1,-1);
\filldraw[fill=white,draw=black!15,line width=0.08pt] (5,-8) rectangle ++(1,-1);
\filldraw[fill=BaselineBlue!31,draw=BaselineBlue!45!black,line width=0.08pt] (6,-8) rectangle ++(1,-1);
\filldraw[fill=InterceptOrange!31,draw=InterceptOrange!55!black,line width=0.08pt] (7,-8) rectangle ++(1,-1);
\filldraw[fill=BaselineBlue!44,draw=BaselineBlue!45!black,line width=0.08pt] (8,-8) rectangle ++(1,-1);
\filldraw[fill=white,draw=black!15,line width=0.08pt] (9,-8) rectangle ++(1,-1);
\filldraw[fill=BaselineBlue!44,draw=BaselineBlue!45!black,line width=0.08pt] (10,-8) rectangle ++(1,-1);
\filldraw[fill=white,draw=black!15,line width=0.08pt] (11,-8) rectangle ++(1,-1);
\filldraw[fill=BaselineBlue!44,draw=BaselineBlue!45!black,line width=0.08pt] (12,-8) rectangle ++(1,-1);
\filldraw[fill=white,draw=black!15,line width=0.08pt] (13,-8) rectangle ++(1,-1);
\filldraw[fill=white,draw=black!15,line width=0.08pt] (14,-8) rectangle ++(1,-1);
\filldraw[fill=BaselineBlue!44,draw=BaselineBlue!45!black,line width=0.08pt] (15,-8) rectangle ++(1,-1);
\filldraw[fill=white,draw=black!15,line width=0.08pt] (0,-9) rectangle ++(1,-1);
\filldraw[fill=white,draw=black!15,line width=0.08pt] (1,-9) rectangle ++(1,-1);
\filldraw[fill=BaselineBlue!31,draw=BaselineBlue!45!black,line width=0.08pt] (2,-9) rectangle ++(1,-1);
\filldraw[fill=InterceptOrange!31,draw=InterceptOrange!55!black,line width=0.08pt] (3,-9) rectangle ++(1,-1);
\filldraw[fill=white,draw=black!15,line width=0.08pt] (4,-9) rectangle ++(1,-1);
\filldraw[fill=BaselineBlue!62,draw=BaselineBlue!45!black,line width=0.08pt] (5,-9) rectangle ++(1,-1);
\filldraw[fill=BaselineBlue!31,draw=BaselineBlue!45!black,line width=0.08pt] (6,-9) rectangle ++(1,-1);
\filldraw[fill=BaselineBlue!31,draw=BaselineBlue!45!black,line width=0.08pt] (7,-9) rectangle ++(1,-1);
\filldraw[fill=white,draw=black!15,line width=0.08pt] (8,-9) rectangle ++(1,-1);
\filldraw[fill=BaselineBlue!44,draw=BaselineBlue!45!black,line width=0.08pt] (9,-9) rectangle ++(1,-1);
\filldraw[fill=BaselineBlue!44,draw=BaselineBlue!45!black,line width=0.08pt] (10,-9) rectangle ++(1,-1);
\filldraw[fill=white,draw=black!15,line width=0.08pt] (11,-9) rectangle ++(1,-1);
\filldraw[fill=white,draw=black!15,line width=0.08pt] (12,-9) rectangle ++(1,-1);
\filldraw[fill=InterceptOrange!44,draw=InterceptOrange!55!black,line width=0.08pt] (13,-9) rectangle ++(1,-1);
\filldraw[fill=white,draw=black!15,line width=0.08pt] (14,-9) rectangle ++(1,-1);
\filldraw[fill=InterceptOrange!44,draw=InterceptOrange!55!black,line width=0.08pt] (15,-9) rectangle ++(1,-1);
\filldraw[fill=BaselineBlue!62,draw=BaselineBlue!45!black,line width=0.08pt] (0,-10) rectangle ++(1,-1);
\filldraw[fill=white,draw=black!15,line width=0.08pt] (1,-10) rectangle ++(1,-1);
\filldraw[fill=BaselineBlue!62,draw=BaselineBlue!45!black,line width=0.08pt] (2,-10) rectangle ++(1,-1);
\filldraw[fill=white,draw=black!15,line width=0.08pt] (3,-10) rectangle ++(1,-1);
\filldraw[fill=white,draw=black!15,line width=0.08pt] (4,-10) rectangle ++(1,-1);
\filldraw[fill=BaselineBlue!62,draw=BaselineBlue!45!black,line width=0.08pt] (5,-10) rectangle ++(1,-1);
\filldraw[fill=BaselineBlue!62,draw=BaselineBlue!45!black,line width=0.08pt] (6,-10) rectangle ++(1,-1);
\filldraw[fill=white,draw=black!15,line width=0.08pt] (7,-10) rectangle ++(1,-1);
\filldraw[fill=BaselineBlue!44,draw=BaselineBlue!45!black,line width=0.08pt] (8,-10) rectangle ++(1,-1);
\filldraw[fill=BaselineBlue!44,draw=BaselineBlue!45!black,line width=0.08pt] (9,-10) rectangle ++(1,-1);
\filldraw[fill=BaselineBlue!88,draw=BaselineBlue!45!black,line width=0.08pt] (10,-10) rectangle ++(1,-1);
\filldraw[fill=white,draw=black!15,line width=0.08pt] (11,-10) rectangle ++(1,-1);
\filldraw[fill=BaselineBlue!44,draw=BaselineBlue!45!black,line width=0.08pt] (12,-10) rectangle ++(1,-1);
\filldraw[fill=InterceptOrange!44,draw=InterceptOrange!55!black,line width=0.08pt] (13,-10) rectangle ++(1,-1);
\filldraw[fill=white,draw=black!15,line width=0.08pt] (14,-10) rectangle ++(1,-1);
\filldraw[fill=white,draw=black!15,line width=0.08pt] (15,-10) rectangle ++(1,-1);
\filldraw[fill=white,draw=black!15,line width=0.08pt] (0,-11) rectangle ++(1,-1);
\filldraw[fill=white,draw=black!15,line width=0.08pt] (1,-11) rectangle ++(1,-1);
\filldraw[fill=white,draw=black!15,line width=0.08pt] (2,-11) rectangle ++(1,-1);
\filldraw[fill=white,draw=black!15,line width=0.08pt] (3,-11) rectangle ++(1,-1);
\filldraw[fill=white,draw=black!15,line width=0.08pt] (4,-11) rectangle ++(1,-1);
\filldraw[fill=white,draw=black!15,line width=0.08pt] (5,-11) rectangle ++(1,-1);
\filldraw[fill=white,draw=black!15,line width=0.08pt] (6,-11) rectangle ++(1,-1);
\filldraw[fill=white,draw=black!15,line width=0.08pt] (7,-11) rectangle ++(1,-1);
\filldraw[fill=white,draw=black!15,line width=0.08pt] (8,-11) rectangle ++(1,-1);
\filldraw[fill=white,draw=black!15,line width=0.08pt] (9,-11) rectangle ++(1,-1);
\filldraw[fill=white,draw=black!15,line width=0.08pt] (10,-11) rectangle ++(1,-1);
\filldraw[fill=white,draw=black!15,line width=0.08pt] (11,-11) rectangle ++(1,-1);
\filldraw[fill=white,draw=black!15,line width=0.08pt] (12,-11) rectangle ++(1,-1);
\filldraw[fill=white,draw=black!15,line width=0.08pt] (13,-11) rectangle ++(1,-1);
\filldraw[fill=white,draw=black!15,line width=0.08pt] (14,-11) rectangle ++(1,-1);
\filldraw[fill=white,draw=black!15,line width=0.08pt] (15,-11) rectangle ++(1,-1);
\filldraw[fill=BaselineBlue!62,draw=BaselineBlue!45!black,line width=0.08pt] (0,-12) rectangle ++(1,-1);
\filldraw[fill=white,draw=black!15,line width=0.08pt] (1,-12) rectangle ++(1,-1);
\filldraw[fill=BaselineBlue!31,draw=BaselineBlue!45!black,line width=0.08pt] (2,-12) rectangle ++(1,-1);
\filldraw[fill=BaselineBlue!31,draw=BaselineBlue!45!black,line width=0.08pt] (3,-12) rectangle ++(1,-1);
\filldraw[fill=white,draw=black!15,line width=0.08pt] (4,-12) rectangle ++(1,-1);
\filldraw[fill=white,draw=black!15,line width=0.08pt] (5,-12) rectangle ++(1,-1);
\filldraw[fill=BaselineBlue!31,draw=BaselineBlue!45!black,line width=0.08pt] (6,-12) rectangle ++(1,-1);
\filldraw[fill=InterceptOrange!31,draw=InterceptOrange!55!black,line width=0.08pt] (7,-12) rectangle ++(1,-1);
\filldraw[fill=BaselineBlue!44,draw=BaselineBlue!45!black,line width=0.08pt] (8,-12) rectangle ++(1,-1);
\filldraw[fill=white,draw=black!15,line width=0.08pt] (9,-12) rectangle ++(1,-1);
\filldraw[fill=BaselineBlue!44,draw=BaselineBlue!45!black,line width=0.08pt] (10,-12) rectangle ++(1,-1);
\filldraw[fill=white,draw=black!15,line width=0.08pt] (11,-12) rectangle ++(1,-1);
\filldraw[fill=BaselineBlue!44,draw=BaselineBlue!45!black,line width=0.08pt] (12,-12) rectangle ++(1,-1);
\filldraw[fill=white,draw=black!15,line width=0.08pt] (13,-12) rectangle ++(1,-1);
\filldraw[fill=white,draw=black!15,line width=0.08pt] (14,-12) rectangle ++(1,-1);
\filldraw[fill=BaselineBlue!44,draw=BaselineBlue!45!black,line width=0.08pt] (15,-12) rectangle ++(1,-1);
\filldraw[fill=white,draw=black!15,line width=0.08pt] (0,-13) rectangle ++(1,-1);
\filldraw[fill=white,draw=black!15,line width=0.08pt] (1,-13) rectangle ++(1,-1);
\filldraw[fill=InterceptOrange!31,draw=InterceptOrange!55!black,line width=0.08pt] (2,-13) rectangle ++(1,-1);
\filldraw[fill=BaselineBlue!31,draw=BaselineBlue!45!black,line width=0.08pt] (3,-13) rectangle ++(1,-1);
\filldraw[fill=white,draw=black!15,line width=0.08pt] (4,-13) rectangle ++(1,-1);
\filldraw[fill=InterceptOrange!62,draw=InterceptOrange!55!black,line width=0.08pt] (5,-13) rectangle ++(1,-1);
\filldraw[fill=InterceptOrange!31,draw=InterceptOrange!55!black,line width=0.08pt] (6,-13) rectangle ++(1,-1);
\filldraw[fill=InterceptOrange!31,draw=InterceptOrange!55!black,line width=0.08pt] (7,-13) rectangle ++(1,-1);
\filldraw[fill=white,draw=black!15,line width=0.08pt] (8,-13) rectangle ++(1,-1);
\filldraw[fill=InterceptOrange!44,draw=InterceptOrange!55!black,line width=0.08pt] (9,-13) rectangle ++(1,-1);
\filldraw[fill=InterceptOrange!44,draw=InterceptOrange!55!black,line width=0.08pt] (10,-13) rectangle ++(1,-1);
\filldraw[fill=white,draw=black!15,line width=0.08pt] (11,-13) rectangle ++(1,-1);
\filldraw[fill=white,draw=black!15,line width=0.08pt] (12,-13) rectangle ++(1,-1);
\filldraw[fill=BaselineBlue!44,draw=BaselineBlue!45!black,line width=0.08pt] (13,-13) rectangle ++(1,-1);
\filldraw[fill=white,draw=black!15,line width=0.08pt] (14,-13) rectangle ++(1,-1);
\filldraw[fill=BaselineBlue!44,draw=BaselineBlue!45!black,line width=0.08pt] (15,-13) rectangle ++(1,-1);
\filldraw[fill=white,draw=black!15,line width=0.08pt] (0,-14) rectangle ++(1,-1);
\filldraw[fill=white,draw=black!15,line width=0.08pt] (1,-14) rectangle ++(1,-1);
\filldraw[fill=white,draw=black!15,line width=0.08pt] (2,-14) rectangle ++(1,-1);
\filldraw[fill=white,draw=black!15,line width=0.08pt] (3,-14) rectangle ++(1,-1);
\filldraw[fill=white,draw=black!15,line width=0.08pt] (4,-14) rectangle ++(1,-1);
\filldraw[fill=white,draw=black!15,line width=0.08pt] (5,-14) rectangle ++(1,-1);
\filldraw[fill=white,draw=black!15,line width=0.08pt] (6,-14) rectangle ++(1,-1);
\filldraw[fill=white,draw=black!15,line width=0.08pt] (7,-14) rectangle ++(1,-1);
\filldraw[fill=white,draw=black!15,line width=0.08pt] (8,-14) rectangle ++(1,-1);
\filldraw[fill=white,draw=black!15,line width=0.08pt] (9,-14) rectangle ++(1,-1);
\filldraw[fill=white,draw=black!15,line width=0.08pt] (10,-14) rectangle ++(1,-1);
\filldraw[fill=white,draw=black!15,line width=0.08pt] (11,-14) rectangle ++(1,-1);
\filldraw[fill=white,draw=black!15,line width=0.08pt] (12,-14) rectangle ++(1,-1);
\filldraw[fill=white,draw=black!15,line width=0.08pt] (13,-14) rectangle ++(1,-1);
\filldraw[fill=white,draw=black!15,line width=0.08pt] (14,-14) rectangle ++(1,-1);
\filldraw[fill=white,draw=black!15,line width=0.08pt] (15,-14) rectangle ++(1,-1);
\filldraw[fill=BaselineBlue!62,draw=BaselineBlue!45!black,line width=0.08pt] (0,-15) rectangle ++(1,-1);
\filldraw[fill=white,draw=black!15,line width=0.08pt] (1,-15) rectangle ++(1,-1);
\filldraw[fill=white,draw=black!15,line width=0.08pt] (2,-15) rectangle ++(1,-1);
\filldraw[fill=BaselineBlue!62,draw=BaselineBlue!45!black,line width=0.08pt] (3,-15) rectangle ++(1,-1);
\filldraw[fill=white,draw=black!15,line width=0.08pt] (4,-15) rectangle ++(1,-1);
\filldraw[fill=InterceptOrange!62,draw=InterceptOrange!55!black,line width=0.08pt] (5,-15) rectangle ++(1,-1);
\filldraw[fill=white,draw=black!15,line width=0.08pt] (6,-15) rectangle ++(1,-1);
\filldraw[fill=InterceptOrange!62,draw=InterceptOrange!55!black,line width=0.08pt] (7,-15) rectangle ++(1,-1);
\filldraw[fill=BaselineBlue!44,draw=BaselineBlue!45!black,line width=0.08pt] (8,-15) rectangle ++(1,-1);
\filldraw[fill=InterceptOrange!44,draw=InterceptOrange!55!black,line width=0.08pt] (9,-15) rectangle ++(1,-1);
\filldraw[fill=white,draw=black!15,line width=0.08pt] (10,-15) rectangle ++(1,-1);
\filldraw[fill=white,draw=black!15,line width=0.08pt] (11,-15) rectangle ++(1,-1);
\filldraw[fill=BaselineBlue!44,draw=BaselineBlue!45!black,line width=0.08pt] (12,-15) rectangle ++(1,-1);
\filldraw[fill=BaselineBlue!44,draw=BaselineBlue!45!black,line width=0.08pt] (13,-15) rectangle ++(1,-1);
\filldraw[fill=white,draw=black!15,line width=0.08pt] (14,-15) rectangle ++(1,-1);
\filldraw[fill=BaselineBlue!88,draw=BaselineBlue!45!black,line width=0.08pt] (15,-15) rectangle ++(1,-1);
\node[anchor=north] at (0.5,-16.15) {0};
\node[anchor=east] at (-0.15,-0.5) {0};
\node[anchor=north] at (4.5,-16.15) {4};
\node[anchor=east] at (-0.15,-4.5) {4};
\node[anchor=north] at (8.5,-16.15) {8};
\node[anchor=east] at (-0.15,-8.5) {8};
\node[anchor=north] at (12.5,-16.15) {12};
\node[anchor=east] at (-0.15,-12.5) {12};
\node[anchor=north] at (15.5,-16.15) {15};
\node[anchor=east] at (-0.15,-15.5) {15};
\node[anchor=north] at (8,-17.15) {column};
\node[rotate=90] at (-1.15,-8) {row};
\end{scope}
\end{tikzpicture}
\caption{Real parts of the reconstructed $16\times16$ reduced density
matrices. Blue denotes positive entries and orange denotes negative entries.
The baseline is pure. Both modified circuits have purity $1/2$. The advanced
circuit preserves the baseline diagonal but changes the coherence pattern away
from the diagonal.}
\label{fig:caseii-density-matrices}
\end{figure}

There is a direct way to see where that lost purity went. The branch-resolved
amplitudes are plotted in \cref{fig:caseii-branch-amplitudes}. In the baseline,
the second fifth-wire branch is empty. In both modified circuits, two branches
remain populated.

\begin{figure}[tbp]
\centering
\begin{tikzpicture}
\begin{groupplot}[
group style={group size=1 by 3,vertical sep=0.82cm},
width=0.92\textwidth,height=0.185\textwidth,
xmin=-0.7,xmax=15.7,ymin=-0.39,ymax=0.39,
ytick={-0.353553,0,0.353553},
yticklabels={$-1/(2\sqrt2)$,$0$,$1/(2\sqrt2)$},
ylabel={real amplitude},
axis line style={black!65},tick style={black!65},
title style={font=\small\bfseries},
grid=major,grid style={black!10},
]
\nextgroupplot[title={Baseline},xtick=\empty]
\addplot[ybar,bar width=2.2pt,fill=BaselineBlue!82,draw=BaselineBlue!95!black] coordinates {(-0.14,0.35355339374) (0.86,0) (1.8599999999999999,0.249999997775) (2.86,0.249999997775) (3.86,0) (4.86,0.35355339374) (5.86,0.249999997775) (6.86,-0.249999997775) (7.86,0.249999997775) (8.86,0.249999997775) (9.86,0.35355339374) (10.86,0) (11.86,0.249999997775) (12.86,-0.249999997775) (13.86,0) (14.86,0.35355339374)};
\addplot[ybar,bar width=2.2pt,fill=BaselineBlue!22,draw=BaselineBlue!95!black,postaction={pattern=north east lines,pattern color=BaselineBlue!80!black}] coordinates {(0.14,0) (1.1400000000000001,0) (2.14,0) (3.14,0) (4.14,0) (5.14,0) (6.14,0) (7.14,0) (8.14,0) (9.14,0) (10.14,0) (11.14,0) (12.14,0) (13.14,0) (14.14,0) (15.14,0)};
\nextgroupplot[title={Intercept--resend},xtick=\empty]
\addplot[ybar,bar width=2.2pt,fill=InterceptOrange!82,draw=InterceptOrange!95!black] coordinates {(-0.14,0.17677669687) (0.86,0.17677669687) (1.8599999999999999,0.249999997775) (2.86,0) (3.86,0.17677669687) (4.86,0.17677669687) (5.86,0.249999997775) (6.86,0) (7.86,0.249999997775) (8.86,0.249999997775) (9.86,0.35355339374) (10.86,0) (11.86,0) (12.86,0) (13.86,0) (14.86,0)};
\addplot[ybar,bar width=2.2pt,fill=InterceptOrange!22,draw=InterceptOrange!95!black,postaction={pattern=north east lines,pattern color=InterceptOrange!80!black}] coordinates {(0.14,0.17677669687) (1.1400000000000001,-0.17677669687) (2.14,0) (3.14,0.249999997775) (4.14,0.17677669687) (5.14,-0.17677669687) (6.14,0) (7.14,0.249999997775) (8.14,0.249999997775) (9.14,-0.249999997775) (10.14,0) (11.14,0.35355339374) (12.14,0) (13.14,0) (14.14,0) (15.14,0)};
\nextgroupplot[title={Advanced source access},xtick={0,1,...,15},xticklabels={$0000$,$0001$,$0010$,$0011$,$0100$,$0101$,$0110$,$0111$,$1000$,$1001$,$1010$,$1011$,$1100$,$1101$,$1110$,$1111$},x tick label style={rotate=55,anchor=east,font=\scriptsize},xlabel={display-register basis state}]
\addplot[ybar,bar width=2.2pt,fill=AdvancedTeal!82,draw=AdvancedTeal!95!black] coordinates {(-0.14,0.249999997775) (0.86,0) (1.8599999999999999,0.249999997775) (2.86,0) (3.86,0) (4.86,0.249999997775) (5.86,0.249999997775) (6.86,0) (7.86,0.17677669687) (8.86,0.17677669687) (9.86,0.35355339374) (10.86,0) (11.86,0.17677669687) (12.86,-0.17677669687) (13.86,0) (14.86,0)};
\addplot[ybar,bar width=2.2pt,fill=AdvancedTeal!22,draw=AdvancedTeal!95!black,postaction={pattern=north east lines,pattern color=AdvancedTeal!80!black}] coordinates {(0.14,0.249999997775) (1.1400000000000001,0) (2.14,0) (3.14,0.249999997775) (4.14,0) (5.14,-0.249999997775) (6.14,0) (7.14,-0.249999997775) (8.14,0.17677669687) (9.14,-0.17677669687) (10.14,0) (11.14,0) (12.14,0.17677669687) (13.14,0.17677669687) (14.14,0) (15.14,0.35355339374)};
\end{groupplot}
\end{tikzpicture}
\caption{Real amplitudes separated by the fifth exported wire. Solid bars show
the $e=0$ branch and hatched bars show the $e=1$ branch. The baseline occupies
one branch. The modified circuits distribute amplitude across two orthogonal
branches.}
\label{fig:caseii-branch-amplitudes}
\end{figure}

For each modified circuit, the complete five-wire state admits the form
\begin{equation}
  \ket{\Psi_x}
  =\frac{1}{\sqrt2}
   \left(
     \ket{\chi^{(x)}_0}\ket{0}_e
     +\ket{\chi^{(x)}_1}\ket{1}_e
   \right),
  \qquad
  \braket{\chi^{(x)}_0|\chi^{(x)}_1}=0.
  \label{eq:caseii-schmidt}
\end{equation}
Tracing out the fifth wire gives
\begin{equation}
  \rho_x
  =\frac12\ket{\chi^{(x)}_0}\!\bra{\chi^{(x)}_0}
  +\frac12\ket{\chi^{(x)}_1}\!\bra{\chi^{(x)}_1},
\end{equation}
so
\begin{equation}
  \Tr(\rho_x^2)=\frac12,
  \qquad
  S(\rho_x)=1\ \text{bit}.
  \label{eq:caseii-one-bit}
\end{equation}

At first glance, the advanced circuit appears to put everything back. The exact
calculation says something more precise: it puts back the entire
computational-basis histogram while leaving one bit of branch correlation in a
wire the histogram does not include. It restores the visible answer to one
question, not the original state.

\begin{rem}[What the exact data establishes]
The advanced export is perfectly stealthy against the supplied four-wire
$Z$-basis histogram test in the ideal simulator output. It is not stealthy
against complete state verification, complementary-basis tests, or an optimal
distinguishing measurement.
\end{rem}

\subsection{Fixed-input fifth-wire identification}

The randomized export establishes hidden correlation but averages together the
preparation labels and measurement branches. The accompanying exact simulator
therefore fixes Alice's value $v$, Alice's basis $b$, and the spy's measurement
basis $e$, then enumerates the measurement result $r_E$ rather than sampling
shots. Let $\rho_4^{(v,b,e,r_E)}$ denote the final state of the retained physical
wire $q_4$ after routing.

\begin{proposition}[Branch-conditioned retained-wire state]
In the reconstructed source-access model,
\begin{equation}
  \rho_4^{(v,b,e,r_E)}=\ket{r_E}\!\bra{r_E}.
  \label{eq:caseii-branch-retained-state}
\end{equation}
After averaging over the spy's measurement outcome,
\begin{equation}
  \rho_4^{(v,b,e)}=
  \begin{cases}
    \ket{v}\!\bra{v}, & e=b,\\[2pt]
    I/2, & e\neq b.
  \end{cases}
  \label{eq:caseii-fixed-input-theorem}
\end{equation}
\end{proposition}

\begin{proof}
When $e=b$, measurement occurs in the preparation basis and returns
$r_E=v$ deterministically. When $e\neq b$, the BB84 state is unbiased in the
measurement basis, so $r_E$ is uniform. The SWAP permutation in
\cref{eq:caseii-semantic-permutation} then moves the post-measurement $q_2$
branch to $q_4$.
\end{proof}

Consequently, for fixed $b$ and $e$,
\begin{equation}
D\!\left(\rho_4^{(0,b,e)},\rho_4^{(1,b,e)}\right)
=\begin{cases}1,&e=b,\\0,&e\neq b.\end{cases}
\label{eq:caseii-value-distinguishability}
\end{equation}
The optimal single-shot value guess is therefore perfect in the matching-basis
case and no better than random in the mismatched-basis case. If $e$ is uniform
and later forgotten, the trace distance is $1/2$, giving a $75\%$ optimal value
guess and $1-h_2(1/4)\approx0.188722$ bits of mutual information. If the spy
retains $e$ as side information, the average value information is $0.5$ bits per
attacked transmission.

The same branch enumeration checks the legitimate subsystem after routing. For
every fixed $(v,b,e,r_E)$ branch, the reconstructed victim state has fidelity
$1$ with its clean routed reference and trace distance at most
$1.23\times10^{-32}$ in the numerical run. Under the provisional acceptance
rule
\begin{equation}
  \mathrm{accept}=(c\neq b)\lor(r_B=v),
  \label{eq:caseii-provisional-acceptance}
\end{equation}
matched-basis acceptance is also unity to floating-point precision.

\begin{rem}[Remaining detector ambiguity]
Equation~\eqref{eq:caseii-provisional-acceptance} is a reasonable BB84-style
sifting rule, but it has not yet been proven to be the exact Boolean acceptance
register implemented by the full original educational circuit. The retained
wire's value leakage is established for the reconstructed source-access model;
the original detector-register semantics remain an explicit roadmap item.
\end{rem}

\subsection{Why this case matters beyond the educational circuit}

The lesson is not that a state government can automatically recover arbitrary
encrypted quantum data. Encryption and no-cloning still matter, and arbitrary
nonorthogonal states cannot be learned perfectly while remaining unchanged.
The narrower and more defensible lesson is that a privileged compiler, runtime,
or source implementation may have access to classical preparation labels,
orthogonal intermediate data, or reusable workspace that a channel-only
security model never grants to an attacker.

That creates a gap between two statements that sound similar but are not:

\begin{align}
  &\text{``The final bitstring distribution is unchanged.''}\label{eq:caseii-same-histogram}\\
  &\text{``The computation and its information flow are unchanged.''}\label{eq:caseii-same-computation}
\end{align}

The advanced export satisfies the first statement for the tested register and
fails the second. That difference is the bridge from this educational example
to the broader malicious-reuse threat model developed later in the paper.
